\section{Rings and Ideals (Done)}

\subsection{Exercise}\label{1.1}

\begin{proof}
By assumption, $x^m = 0$ for some $m$. Then, 
\[ (1 + x) \sum_{i=0}^{m-1} (-1)^i x^i = 1 + (-1)^{m-1}x^m = 1, \] 
so $1 + x$ is a unit.

Let $x$ be nilpotent and $u$ be a unit. Then, $u + x = u (1 + u^{-1} x)$. But $u^{-1} x$ is nilpotent since $x$ is. By the previous result, $1 + u^{-1} x$ is a unit and hence $u (1 + u^{-1} x) = u + x$ is also a unit.
\end{proof}

\subsection{Exercise}\label{1.2}

\begin{proof}[(i) Proof.]
We will prove this by induction. Suppose $f$ has degree $0$ with inverse $g = b_0 + \dots + b_m x^m$. Then,
$1 = fb_0 + \dots + fb_m x^m$. Therefore, since $f \in A$ and $fb_0 = 1$, $f$ is a unit. The converse is trivial since $f=a_0$,
proving the base case.

Now, suppose that for all polynomials of degree $n-1$, the proposition holds. Suppose $f$ has degree $n$ and 
has an inverse $g = b_0 + \dots + b_mx^m$. Then,
\begin{align*}
1 &= \left(\sum_{i=0}^n a_i x^i\right) \cdot \left(\sum_{j=0}^m b_j x^j\right) \\
&= \sum_{k=0}^{n+m}\left(\sum_{i=0}^k a_i b_{k-i}\right)x^i.
\end{align*}
This implies $a_n b_m = 0$. We'll show by induction on $r$ that $a_n^{r+1} b_{m-r} = 0$, the base case already
proven. We assume this is true for all $r < l$. Multiplying by $a_n^l$,
\begin{align*}
a_n^l &= \left(\sum_{i=0}^n a_i x^i\right) \cdot \left(\sum_{j=0}^m a_n^l b_j x^j\right) \\
&= \left(\sum_{i=0}^n a_i x^i\right) \cdot \left(\sum_{j=0}^{m-l} a_n^l b_j x^j\right),
\end{align*}
where we used the inductive hypothesis to remove higher terms: if $j > m-l$, then $l-(m-j+1) \ge 0$, so $a_n^l b_j = a_n^{l-(m-j+1)}(a_n^{m-j+1}b_{m-(m-j)}) = 0$.
The coefficient of highest degree is $a_n \cdot a_n^l b_{m-l}$, so $a_n^{l+1}b_{m-l}=0$. By induction, $a_n^{r+1}b_{m-r}=0$
for all $r$.

From the expansion of $f \cdot g = 1$, we know that $a_0 b_0 = 1$. Then, using our result,
$$a_n^{m+1} = a_0 (a_n^{m+1}b_{m - m}) = a_0 \cdot 0 = 0,$$
so $a_n$ is nilpotent. Then, $- a_n x^n$ is nilpotent. Since $f$ is a unit, by Exercise \ref{1.1},
$f + (-a_n x^n)$ is also a unit. However, it has degree $n-1$. By the inductive hypothesis, $a_0$ is a unit and $a_i$ is nilpotent
for $1 \le i < n$. We've shown $a_n$ is nilpotent, so this direction is complete.

Conversely, suppose $a_0$ is a unit and that $a_i$ is nilpotent for $1 \le i \le n$. Then, $f - a_nx^n$ has degree $n-1$,
so by the inductive hypothesis, $f - a_nx^n$ is a unit since its coefficients have the right properties. 
However, $a_n x^n$ is nilpotent since $a_n$ is, so by Exercise \ref{1.1},
$(f - a_nx^n) + a_n x^n = f$ is a unit, so both directions are complete. By induction, the proposition follows for $f$ of any degree.
\end{proof}

\begin{proof}[(ii) Proof.]
We'll prove this by induction on the degree. The base case is trivial since $f = a_0$ in degree 0. Now, suppose for all degrees 
less than $n$, the proposition holds.

If $f$ has degree $n$ and is nilpotent, there is an $m$ such that $f^m = 0$. Expanding, the top term is $a_n^m x^{nm}$, which is zero. Therefore,
$a_n$ is nilpotent. However, $g = f + (- a_nx^n)$ is also nilpotent since it is the sum of two nilpotent elements (if $a$ and $b$ are the required nilpotent
powers, take $a + b$; this power will annihilate all terms in the binomial expansion of $g$). $g$ has degree at most $n-1$, so by the inductive hypothesis
$a_0,\ldots,a_{n-1}$ are all nilpotent. Since $a_n$ is nilpotent, this direction is complete.

Conversely, if $f$ has degree $n$ and all coefficients are nilpotent, take $f = g + a_n x^n$, where $g$ has degree at most $n-1$. By the inductive hypothesis,
$g$ is nilpotent, so $f$ is the sum of two nilpotent elements. Thus, $f$ is nilpotent. By induction, the proposition follows for $f$ of any degree.
\end{proof}

\begin{proof}[(iii) Proof.]
The reverse implication is by definition. For the forward implication, suppose $f$ is a zero-divisor. The set of all possible $g$ such that
$fg = 0$ has minimal degree $m$ since it is non-empty. Choose $g$ of degree $m$ such that $fg = 0$. If the highest terms are $a_n x^n$ and $b_m x^m$, 
then $a_nb_mx^{n+m} = 0$ by $fg = 0$, so $a_n b_m = 0$. Note that $a_n g$ annihilates the term of degree $m$ and annihilates $f$, so $a_n g = 0$; 
otherwise, we would be violating the minimality of $m$. 

We will show by induction that $a_{n-r}g = 0$, the base case just proven. For the inductive hypothesis, suppose $a_{n-r}g=0$ for all $r < k$. Then,
$$0 = g\sum_{i=0}^na_ix^i = \sum_{i=0}^n(ga_i)x^i = \sum_{i=0}^{n-k}(ga_i)x^i.$$
The top term is $b_m a_{n-k} x^{m+n-k}$, so $b_m a_{n-k} = 0$. Thus, $a_{n-k}g$ has degree less than $m$ and annihilates $f$. 
Again, to preserve the minimality of $m$, we have $a_{n-k} g = 0$. By induction, $a_{n-r}g = 0$ for all $r$.

For any $a_i$, we have $a_ig = 0$. Thus,
$$0 = a_i\left(b_0 + \dots + b_m x^m\right) = a_ib_0 + \dots + a_ib_m x^m.$$
By definition, $b_m \neq 0$, yet $a_ib_m x^m = 0$. So, $a_i b_m = 0$ for all $i$. Thus,
$$b_m f = b_m \sum_{i=0}^n (b_m a_i) x^i = 0,$$
proving the claim.
\end{proof}

\begin{remark}
For use in Exercise \ref{1.3}, we just proved: if $fg = 0$, $g$ annihilates every coefficient of $f$.
\end{remark}

\begin{proof}[(iv) Proof.]
Suppose $fg$ is not primitive. Then, the ideal generated by its coefficients lie in a maximal ideal $\mathfrak{m}$.
Thus, under the ring homomorphism $A[x] \to (A/\mathfrak{m})[x]$, $0 = \overline{fg} = \overline{f}\overline{g}$.
If $\overline{f}$ is a zero divisor, then by (iii), there is a nonzero $a \in A/\mathfrak{m}$ such that
$a\overline{f} = 0$. But, $a$ is a unit since $A/\mathfrak{m}$ is a field, so $\overline{f}=0$. The same argument applies
to $\overline{g}$. Thus, either $\overline{f}$ or $\overline{g}$ is $0$; pulling back whichever is zero, its coefficients must
all lie in $\mathfrak{m}$, violating primitivity.

Conversely, suppose $f$ or $g$ is not primitive. Without loss of generality, we can say $f$ is not primitive, so its coefficients lie in
a maximal ideal $\mathfrak{m}$. Then,
$$fg = \sum_{k=0}^{n+m}\left(\sum_{i=0}^ka_ib_{k-i}\right)x^k.$$
Obviously, $a_i b_{k-i} \in \mathfrak{m}$ since $a_i \in \mathfrak{m}$, so any finite sum of these terms is also. Therefore, the coefficients
of $fg$ all lie in $\mathfrak{m}$, so $fg$ is not primitive, concluding the proof.
\end{proof}

\subsection{Exercise}\label{1.3}

Here, we will frequently use the fact that
$$A[x_1,\ldots,x_n] \cong A[x_1,\ldots,x_{n-1}][x_n],$$
which is easy to see by simply grouping terms on the left-hand side by $x_n$-degree.

\begin{proof}[(ii) Proof.]
We'll prove this by induction, using Exercise \ref{1.2}(ii), which establishes the base case. Now, assume the inductive hypothesis up to $n-1$.

Let $f \in A[x_1,\ldots,x_n]$. We can write $f = f_0 + \cdots + f_m x_n^m \in A[x_1,\ldots,x_{n-1}][x_n]$. Assume $f$ is nilpotent. By Exercise \ref{1.2}(ii),
each $f_i$ is nilpotent. Then, by the inductive hypothesis for (ii), all coefficients are nilpotent for each $f_i$ and thus for $f$. Conversely, if all
coefficients of $f$ are nilpotent, then all coefficients of each $f_i$ are nilpotent. By the inductive hypothesis for (ii), each $f_i$ is nilpotent. By
Exercise \ref{1.2}(ii), $f$ is nilpotent. By induction, the result follows.
\end{proof}

\begin{proof}[(i) Proof.]
We'll prove this by induction, using Exercise \ref{1.2}(i), which establishes the base case. Now, assume the inductive hypothesis up to $n-1$.

Suppose $f$ is a unit in $A[x_1,\ldots,x_n] = A[x_1,\ldots,x_{n-1}][x_n]$. Then, $f = f_0 + f_1 x_n + \cdots + f_m x_n^m$, where 
$f_i \in A[x_1,\ldots,x_{n-1}]$. By the inductive hypothesis for (i), $f_0$ is a unit and $f_i$ is nilpotent for $i > 0$. By the inductive hypothesis 
for (ii), all the coefficients of $f_i$ are nilpotent for $i > 0$. By the inductive hypothesis for $f_0$, the constant coefficient is a unit, and all
other coefficients are nilpotent. Thus, viewing $f$ in $A[x_1,\ldots,x_n]$, the constant term is a unit and all other coefficients are nilpotent.

Conversely, suppose $f$ in $A[x_1,\ldots,x_n]$ with the constant term a unit and all other coefficients nilpotent. Then, we can group terms so that
$f = f_0 + f_1 x_n + \cdots + f_m x_n^m$, viewing $f \in A[x_1,\ldots,x_{n-1}][x_n]$. Then, $f_0$ has a unit for a constant term, with all other 
coefficients being nilpotent. Thus, $f_0$ is a unit by the inductive hypothesis for (i). Likewise, $f_i$ for $i > 0$ has all nilpotent coefficients, so
Exercise \ref{1.3}(ii) shows that $f_i$ is nilpotent. By Exercise \ref{1.2}(i), $f$ is a unit. By induction, (i) is proven in the general case.
\end{proof}

\begin{proof}[(iii) Proof.]
First, we start with a lemma: we want to show by induction that for $f,g \in A[x_1,\ldots,x_n]$, $fg = 0$, we have that $g$ annihilates every 
coefficient of $f$. The base case is shown in Exercise \ref{1.2}(iii).

Suppose $f$ is a zero-divisor and $f \in A[x_1,\ldots,x_n] \cong A[x_1,\ldots,x_{n-1}][x_n]$. The set of all possible $g$ such that $fg = 0$ has minimal 
degree $m$ in $x_n$ since it is non-empty. Choose $g$ of degree $m$ in $x_n$ such that $fg = 0$. If the highest terms are $f_k x_n^k$ and $g_l x_n^l$, 
then $f_k g_l x^{k+l} = 0$ by $fg = 0$, so $f_k g_l = 0$. Note that $f_k g$ annihilates the term of degree $m$ and annihilates $f$, so $f_k g = 0$; 
otherwise, we would be violating the minimality of $m$.

We will show by induction that $f_{k-r}g = 0$, the base case just proven. For the inductive hypothesis, suppose $f_{k-r}g=0$ for all $r < q$. Then,
$$0 = g\sum_{i=0}^k f_i x_n^i = \sum_{i=0}^k (gf_i)x_n^i = \sum_{i=0}^{k-q} (gf_i)x_n^i.$$
The top term is $g_l f_{k-q} x_n^{l+k-q}$, so $g_l f_{k-q} = 0$. Thus, $f_{k-q} g$ has degree less than $m$ and annihilates $f$. 
Again, to preserve the minimality of $m$, we have $f_{k-q} g = 0$. By induction, $f_{k-r}g = 0$ for all $r$.

For any $f_i$, we have $f_ig = 0$. Thus,
$$0 = f_i\left(g_0 + \dots + g_m x_n^m\right) = f_i g_0 + \dots + f_i g_m x_n^m.$$ 
For any $g_j$, we have $f_i g_j x_n^j = 0$. So, $f_i g_j = 0$ for all $i$ and $j$. By induction,
since $f_i,g_j \in A[x_1,\ldots,x_{n_1}]$, $g_j$ annihilates every coefficient of $f_i$. Since this
is true for all $j$, $g$ annihilates every coefficient of $f_i$ for all $i$. By induction, for all $n$,
if $f,g \in A[x_1,\ldots,x_n]$ and $fg = 0$, then $g$ annihilates every coefficient of $f$ individually. This proves the lemma.

With this lemma, we can prove (iii). The reverse implication is by definition. For the forward implication, we prove (iii) 
by induction, using the base case established by Exercise \ref{1.2}(iii). Assume the inductive hypothesis up to $n-1$ and that $f,g \in A[x_1,\ldots,x_n]$, 
$fg = 0$. By our lemma, and since $f$ and $g$ are symmetrically defined, $f$ annihilates every coefficient of $g$ individually; since $g$ is nonzero,
it has some nonzero coefficient $g_{\beta} \in A$. Thus, $g_{\beta} f = 0$. By induction, if for any $n$, $f,g \in A[x_1,\ldots,x_n]$ and $fg = 0$,
then $af = 0$ for some nonzero $a \in A$.
\end{proof}

\begin{proof}[(iv) Proof.]
Suppose $fg$ is not primitive. Then, the ideal generated by its coefficients lie in a maximal ideal $\mathfrak{m}$.
Thus, under the ring homomorphism $A[x_1,\ldots,x_n] \to (A/\mathfrak{m})[x_1,\ldots,x_n]$, $0 = \overline{fg} = \overline{f}\cdot\overline{g}$.
If $\overline{f}$ is a zero divisor, then by (iii), there is a nonzero $a \in A/\mathfrak{m}$ such that
$a\overline{f} = 0$. But, $a$ is a unit since $A/\mathfrak{m}$ is a field, so $\overline{f}=0$. The same argument applies
to $\overline{g}$. Thus, either $\overline{f}$ or $\overline{g}$ is $0$; pulling back whichever is zero, its coefficients must
all lie in $\mathfrak{m}$, violating primitivity.

Conversely, suppose $f$ or $g$ is not primitive. Without loss of generality, we can say $f$ is not primitive, so its coefficients lie in
a maximal ideal $\mathfrak{m}$. Then, consider the same map as before, $A[x_1,\ldots,x_n] \to (A/\mathfrak{m})[x_1,\ldots,x_n]$. Then,
$\overline{fg} = \overline{f}\cdot\overline{g} = 0 \cdot \overline{g} = 0$. In particular, all the coefficients of $fg$ lie in $\mathfrak{m}$, 
so $fg$ is not primitive, concluding the proof.
\end{proof}

\subsection{Exercise}\label{1.4}

\begin{proof}
Since all maximal ideals are prime ideals, the nilradical (intersection of all prime ideals) is contained in the Jacobson radical 
(intersection of all maximal ideals) for any ring. Thus, we need to show that any $f(x)$ in the Jacobson radical is nilpotent.

By Proposition 1.9, if $f(x)$ is in the Jacobson radical, then $1 + xf(x)$ is a unit in $A[x]$. We can write this as 
$$1 + xf(x) = 1 + a_0x + \cdots + a_nx^{n+1}.$$
By Exercise \ref{1.2}(i), $a_0,\ldots,a_n$ are nilpotent. Thus, by Exercise \ref{1.2}(ii), $f$ is nilpotent.
\end{proof}

\subsection{Exercise}\label{1.5}

\begin{proof}[(i) Proof.]
Suppose $f = \sum_i a_i x^i \in A[[x]]$ is a unit so that there is some $g = \sum_i b_ix^i \in A[[x]]$ yielding $fg = 1$. 
Expanding, the constant term is $a_0 b_0 = 1$. Thus, $a_0$ is a unit in $A$. 

We will prove the converse by induction; we will show that for all $n$, we can construct some $g_n \in A[x]$ (realized as a subring of $A[[x]]$) of degree $n-1$ 
such that $fg_n - 1 \in (x^n)$ and $g_n \equiv g_r \mod (x^r)$ for $r < n$. 
Suppose $f = \sum_i a_i x^i$ and $a_0$ is a unit so that there is a $b_0$ yielding $a_0 b_0 = 1$.
Let $g_0 = b_0$. Then, 
$$fg_0 - 1 = \sum_{i=0}^{\infty} a_ib_0 x^i - 1 = \sum_{i=1}^{\infty} a_ib_0 x^i + (a_0b_0 - 1) = \sum_{i=1}^{\infty} a_ib_0 x^i \in (x^1),$$
the ideal generated by $x$. This proves the base case. Our inductive hypothesis is that there is some $g_{n-1} = \sum_{i=0}^{n-1} b_ix^i$ such that 
$fg_{n-1} - 1 \equiv 0 \mod (x^n)$. Define $b_n = - a_0^{-1}\sum_{i=1}^n a_i b_{n-i}$, recalling that $a_0$ is a unit. We claim that $g_n = \sum_{i=0}^n b_i x^i$
satisfies the proper congruence. We have
\begin{align*}
fg_n &= \sum_{k=0}^{\infty} \left(\sum_{i=0}^k a_i b_{k-i} \right) x^k \\
&= 1 + \left(\sum_{i=0}^n a_i b_{n-i} \right)x^n + \sum_{k=n+1}^{\infty} \left(\sum_{i=0}^k a_i b_{k-i} \right) x^k \\
&\equiv 1 + \left(a_0b_n + \sum_{i=1}^n a_i b_{n-i} \right)x^n \mod (x^{n+1}) \\
&\equiv 1 + \left(a_0\left(- a_0^{-1}\sum_{i=1}^n a_i b_{n-i}\right) + \sum_{i=1}^n a_i b_{n-i} \right)x^n \mod (x^{n+1}) \\
&\equiv 1 + 0 \cdot x^n \mod (x^{n+1}) \\
&\equiv 1 \mod (x^{n+1}). \\
\end{align*}
Note that since only the $x^n$ coefficient was changed, the compatibility requirement holds: i.e., $g_n \equiv g_r \mod (x^r)$ for $r < n$. Thus, if we let
$g = \sum_{i=0}^{\infty} b_i x^i$, then $fg - 1 \equiv 0 \mod (x^n)$ for all $n$. Once we show that $\bigcap_n (x^n) = (0)$, then we'll have $fg = 1$, 
proving the claim. But this is clear since any nonzero element $h \in A[[x]]$ has some nonzero coefficient $h_j$, so $h \not\in (x^{j+1})$.
\end{proof}

\begin{proof}[(ii) Proof.]
Suppose $f = \sum_i a_i x^i$ is nilpotent. We will show that each $a_i$ is nilpotent by induction. For the base case, note that if $f^m = 0$, then
the constant term is $a_0^m = 0$, so $f_0 = a_0$ is nilpotent. For the inductive step, suppose that $a_i$ is nilpotent for $i < n$. Then, since the 
nilradical is an ideal, $f - \sum_{i=0}^{n-1} a_i x^i = \sum_{i=n}^{\infty} a_i x^i$ is nilpotent. Thus, there is some $m$ such that
$$0 = \left(\sum_{i=n}^{\infty} a_i x^i\right)^m \equiv a_n^m x^{nm} \mod (x^{nm+1}),$$
so $a_n$ is nilpotent. The result follows by induction.
\end{proof}

\begin{remark}
Consider the direct product ring $A = \Pi_{i=2}^{\infty}(\Z / 2^i \Z)$. Then, construct 
$$f = \sum_i a_i x^i,$$
where $a_i = (2,\ldots,2,0,\ldots)$; $2$ appears until slot $i$, whereafter it is $0$ indefinitely. Let $a_0 = a_1 = 0$ since there is no map into the ring.
So, the minimal exponent that annihilates $a_i$ is $i$ whenever $i > 1$. Each $a_i$ is therefore nilpotent.

Note that we have principal ideals $\mathfrak{q}_n = (1,\ldots,1,0,\ldots)$ that project the ring onto $\Pi_{i=2}^n(\Z / 2^i \Z)$; $a_i \in \mathfrak{q}_i$.

On the other hand, for any $m$, consider $f^m$. Then, consider $a_{m(m+1)}$. Note that the coefficient summands are the ordered partitions of 
$m(m+1)$ that have size $m$. There is precisely one summand that is not in the ideal $\mathfrak{q}_m$ and that is $a_{m+1}^m$, so the projection of the 
$x^{m(m+1)}$ coefficient onto $A/\mathfrak{q}_m$ is $a_{m+1}^m$. But, this is nonzero since $a_{m+1}^m$ projects onto $\Z / 2^{m+1} \Z$ as $2^m$. Thus,
$f^m$ has a nonzero $x^{m(m+1)}$ coefficient.

In particular, for general commutative rings with unit (not necessarily Noetherian), if $a_i$ is nilpotent for all $i$, this does not imply that
$f = \sum_{i=0}^{\infty}a_i x^i$ is nilpotent.
\end{remark}

\begin{proof}[(iii) Proof.]
Consider the ring homomorphism $\varphi : A[[x]] \to A$ that sends the power series to its constant term.

For every maximal ideal $\mathfrak{m}$ of $A$, $\varphi^{-1}(\mathfrak{m}) = \mathfrak{m} + xA[[x]]$ is an ideal, the contraction of $\mathfrak{m}$. 
Each of these are contained in maximal ideals of $A[[x]]$, so 
\begin{align*}
\bigcap_{\mathfrak{m}}(\mathfrak{m} + xA[[x]]) &\subset \mathfrak{R}(A[[x]]) \\
\bigcap_{\mathfrak{m}}\mathfrak{m} + xA[[x]] &\subset \mathfrak{R}(A[[x]])\\
\mathfrak{R}(A) + xA[[x]] &\subset \mathfrak{R}(A[[x]])\\
\end{align*}
In particular, if $a_0 \in \mathfrak{R}(A)$ and $f$ has constant term $a_0$, then $f \in \mathfrak{R}(A[[x]])$.

Conversely, assume $a_0$ is not in the Jacobson radical of $A$. Then, there is some maximal ideal $\mathfrak{m}$ in $A$ that does not contain $a_0$. 
But, $\mathfrak{m} + xA[[x]]$ is a non-trivial ideal of $A[[x]]$ under the contraction $\varphi^{-1}(\mathfrak{m})$. If $\mathfrak{m} + xA[[x]]$ is maximal, 
we'd have that $f$ is not in the maximal ideal $\mathfrak{m} + xA[[x]]$ (since its constant term $a_0 \not\in \mathfrak{m}$), so 
$f \not\in \mathfrak{R}(A[[x]])$.

So, the proof is complete once we show that $\mathfrak{m} + xA[[x]]$ is maximal. Let $\mathfrak{r}$ be the maximal ideal containing $\mathfrak{m} + xA[[x]]$. In order for 
these ideals to be unequal, there is some $g \in \mathfrak{r}$ such that $g_0 \not\in \mathfrak{m}$. However, projecting to $A$ (inverse map of homomorphism $A \xhookrightarrow{} A[[x]]$), 
we have an ideal $\mathfrak{r}'$ containing $\mathfrak{m}$ 
properly, so by the maximality of $\mathfrak{m}$, $\mathfrak{r} = A$. Thus, $\mathfrak{r} = A + xA[[x]] = A[[x]]$, violating maximality, so $\mathfrak{m} + A[[x]]$
is maximal.
\end{proof}

\begin{proof}[(iv) Proof.]
First, let's show that for any maximal ideal $\mathfrak{m}$ of $A$, $\mathfrak{m} + xA[[x]]$ is maximal. Let $\mathfrak{r}$ be the maximal ideal containing $\mathfrak{m} + xA[[x]]$. 
In order for these ideals to be unequal, there must be some $g \in \mathfrak{r}$ such that $g_0 \not\in \mathfrak{m}$. However, projecting to $A$ (inverse map of homomorphism $A \xhookrightarrow{} A[[x]]$), 
we have an ideal $\mathfrak{r}'$ containing $\mathfrak{m}$ properly, so by the maximality of $\mathfrak{m}$, $\mathfrak{r} = A$. Thus, $\mathfrak{r} = A + xA[[x]] = A[[x]]$, violating maximality, so 
$\mathfrak{m} + A[[x]]$ is maximal.


Suppose $\mathfrak{m}$ is a maximal ideal of $A[[x]]$. Then, under the ring homomorphism $\varphi : A \xhookrightarrow{} A[[x]]$, $\mathfrak{m}^c = \varphi^{-1}(\mathfrak{m})$
is an ideal of $A$. First of all, if $\mathfrak{m}^c = A$, then $\mathfrak{m}$ contains some series with constant term $1$; by (i), it contains a unit, so $\mathfrak{m} = A[[x]]$,
violating maximality. 

If $\mathfrak{m}^c = 0$, then $\mathfrak{m}$ is contained in the ideal $xA[[x]]$, so $\mathfrak{m} = xA[[x]]$. Note that $xA[[x]]$ is 
maximal when $A$ only has the ideals $0$ and $(1)$ (i.e., $A$ is a field) since for any ideal $I$ of $A$ (not $0$ or $A$), $I + xA[[x]]$ is a 
proper ideal containing $\mathfrak{m}$, violating the maximality of $\mathfrak{m}$. Otherwise, the only ideal containing $xA[[x]]$ is 
$A + xA[[x]] = A[[x]]$, the whole ring. For our purposes, just note that $\mathfrak{m}$ is indeed generated by $x$ and $\mathfrak{m}^c$.

Thus, we are left with the case where $\mathfrak{m}^c$ is neither $0$ nor $A$. Since it is a nontrivial ideal, it is contained in some maximal ideal $\mathfrak{r}$ of $A$. Thus,
$\mathfrak{m} \subset \mathfrak{m}^c + xA[[x]] \subset \mathfrak{r} + xA[[x]]$. By the maximality of $\mathfrak{m}$, $\mathfrak{m}^c = \mathfrak{r}$, and $\mathfrak{m} = \mathfrak{m}^c + xA[[x]]$,
completing the proof.
\end{proof}

\begin{proof}[(v) Proof.]
Let $\mathfrak{p}$ be a prime ideal of $A$. Consider the ring homomorphism $\varphi : A[[x]] \to A$ that maps each series to its constant term. Then,
$\varphi^{-1}(\mathfrak{p}) = \mathfrak{p} + xA[[x]]$ is a prime ideal of $A[[x]]$ since contraction preserves primality. Under the ring homomorphism
$\psi : A \xhookrightarrow{} A[[x]]$, $\mathfrak{p}$ is the contraction of the prime ideal $\mathfrak{p} + xA[[x]]$, proving the claim.
\end{proof}

\subsection{Exercise}\label{1.6}

\begin{proof}
Since all maximal ideals are prime, we know that $\mathfrak{N} \subset \mathfrak{R}$. Suppose to the contrary that the ideal
$\mathfrak{R}$ is not contained in $\mathfrak{N}$. Then, by assumption, there is some nonzero idempotent $e \in \mathfrak{R}$.
Since $e^2 = e$, a simple inductive argument reveals that $e$ cannot be nilpotent. 

Let $\mathfrak{p}$ be any prime ideal. Consider the map $A \to A/\mathfrak{p}$. Because $A/\mathfrak{p}$ is an integral domain and 
$0 = \overline{e^2 - e} = \overline{e}(\overline{e} - 1)$, we have $e \in \mathfrak{p}$ or $e \in 1 + \mathfrak{p}$. If $e \in \mathfrak{p}$
for all prime ideals, then $e$ must also be in nilradical, as it is the intersection of prime ideals. Thus, there is some prime ideal
$\mathfrak{p}'$ such that $e \in 1 + \mathfrak{p}'$.

There is some maximal ideal $\mathfrak{m}'$ containing $\mathfrak{p}'$, so $e \in 1 + \mathfrak{m}'$.
However, since $e \in \mathfrak{R}$, which is the intersection of all maximal ideals, $e \in \mathfrak{m}'$ also.
Together, these two facts imply $1 \in \mathfrak{m}'$, yielding the desired contradiction.
\end{proof}

\subsection{Exercise}\label{1.7}

\begin{proof}
Let $\mathfrak{p}$ be a prime ideal so that $A/\mathfrak{p}$ is an integral domain. For any $y \in A/\mathfrak{p}$
such that $y \neq 0$, there is some $x \in A$, such that $\overline{x} = y$. By assumption, there is an $n > 1$ 
such that $x^n = x$, so $y^n = y$. Note that $y^{n-1} - 1 \in A/\mathfrak{p}$ and 
$$y(y^{n-1} - 1) = y^n - y = 0.$$
Since $A/\mathfrak{p}$ is an integral domain and $y \neq 0$, $y^{n-1} = 1$. In particular, $y$ is a unit. Thus,
all nonzero elements are nonzero, implying that $A/\mathfrak{p}$ is a field and $\mathfrak{p}$ is maximal.
\end{proof}

\subsection{Exercise}\label{1.8}

\begin{proof}
Let $\mathcal{S}$ be the set of prime ideals of $\mathcal{S}$. If $\mathfrak{p} \subset \mathfrak{q}$, we'll define a partial ordering
by $\mathfrak{q} \le \mathfrak{p}$. Now, suppose we have a chain $\mathcal{T} \subset \mathcal{S}$. Let 
$$Q = \bigcap_{\mathfrak{q} \in \mathcal{T}}\mathfrak{q},$$
which, as an intersection of ideals, is itself an ideal.

Suppose $a,b \in A$ and $ab \in Q$. If $a \not\in Q$, then there is some $\mathfrak{r} \in \mathcal{T}$ such that 
$a \not\in \mathfrak{r}$. Since $a \not\in \mathfrak{r}$, for any $\mathfrak{q}$ such that $\mathfrak{q} \subset \mathfrak{r}$, $a\not\in\mathfrak{q}$.
Since $\mathcal{T}$ is a chain,
$$Q = \bigcap_{\mathfrak{q} \le \mathfrak{r}}\mathfrak{q} \cap \bigcap_{\mathfrak{q} \ge \mathfrak{r}}\mathfrak{q}
= \mathfrak{r}\cap \bigcap_{\mathfrak{q} \ge \mathfrak{r}}\mathfrak{q}
= \bigcap_{\mathfrak{q} \ge \mathfrak{r}}\mathfrak{q}.$$
Let $\mathfrak{q} \ge \mathfrak{r}$. Then, since $ab \in Q \subset \mathfrak{q}$ and $\mathfrak{q}$ is prime, we must have $b \in \mathfrak{q}$, as we've
already shown $a$ cannot be in $\mathfrak{r}$. Thus, $b \in \mathfrak{q}$ for all $\mathfrak{q} \ge \mathfrak{r}$, so $b \in Q$. Thus, $Q$ is a prime ideal.

For all $\mathfrak{q} \in \mathcal{T}$, $Q \subset \mathfrak{q}$ ($Q \ge \mathfrak{q}$) and is thus a maximal element. Using Zorn's lemma, $\mathcal{S}$
has at least one maximal element. By our notation, these are minimal elements with respect to inclusion since we used $\ge$ for $\subset$.
\end{proof}

\subsection{Exercise}\label{1.9}

\begin{proof}
If $\mathfrak{a} = \bigcap_{\mathfrak{p}} \mathfrak{p}$ is an intersection of prime ideals, then
$$r(\mathfrak{a}) = r\left(\bigcap_{\mathfrak{p}}\mathfrak{p}\right) = \bigcap_{\mathfrak{p}}r\left(\mathfrak{p}\right) = \bigcap_{\mathfrak{p}}\mathfrak{p} = \mathfrak{a}.$$

Conversely, suppose $r(\mathfrak{a}) = \mathfrak{a}$. By Proposition 1.14, $r(\mathfrak{a}) = \bigcap_{\mathfrak{p}} \mathfrak{p}$, where $\mathfrak{p}$ ranges over the prime ideals
containing $\mathfrak{a}$. Therefore, $\mathfrak{a}$ is an intersection of prime ideals.
\end{proof}

\subsection{Exercise}\label{1.10}

\begin{proof}
(i) $\implies$ (ii). Since $\mathfrak{N}$ is the intersection of all prime ideals, $\mathfrak{N}$ is the unique prime ideal. By Theorem 1.3,
$A$ must have a maximal ideal, but since every maximal ideal is prime and $\mathfrak{N}$ is the only prime ideal, $\mathfrak{N}$ must
be maximal. Hence, it is also the only maximal ideal. By Corollary 1.5, every non-unit is contained in a maximal ideal, so every non-unit
must be in $\mathfrak{N}$, hence nilpotent.

(ii) $\implies$ (iii). Consider the ring homomorphism $A \to A/\mathfrak{N}$. Every element $x\in A$ is either a unit or nilpotent; if $x$ is a unit, 
$\overline{x}$ is a unit, and if $x$ is nilpotent, $x \in \mathfrak{N}$, so $\overline{x} = 0$. Since the projection is surjective, every nonzero element
of the ring $A/\mathfrak{N}$ is a unit, i.e. $A/\mathfrak{N}$ is a field.

(iii) $\implies$ (i). Since $A/\mathfrak{N}$ is a field, $\mathfrak{N}$ a maximal ideal. All maximal ideals are prime.
However, $\mathfrak{N}$ is also the intersection of all prime ideals. Suppose there is another prime ideal $\mathfrak{p}$, not equal to $\mathfrak{N}$.
Thus, there is an element $x \in \mathfrak{p} - \mathfrak{N}$ since $\mathfrak{N} \subset \mathfrak{p}$. Under the map $A \to A/\mathfrak{N}$,
$\mathfrak{x}$ is a unit since it is non-zero. There is a $y$ in $A$ such that $xy = 1 + \mathfrak{N}$. But, $xy \in \mathfrak{p}$ since $x \in \mathfrak{p}$,
so $\mathfrak{p} + \mathfrak{N} = (1)$ and are thus coprime. But $\mathfrak{N} \subset \mathfrak{p}$, so $\mathfrak{p} = (1)$, a contradiction.
\end{proof}

\subsection{Exercise}\label{1.11}

\begin{proof}[(i) Proof.]
For every $x \in A$, 
\begin{align*}
    (1 + x)^2 &= 1 + x \\
    1 + 2x + x^2 &= 1 + x \\
    1 + 2x + x &= 1 + x \\
    2x &= 0.
\end{align*}
\end{proof}

\begin{proof}[(ii) Proof.]
Suppose $x \in A$ represents the class $\bar{x} \in A / \mathfrak{p}$. Then, 
\[ \bar{x}^2 = (\mathfrak{p} + x)(\mathfrak{p} + x) = \mathfrak{p} + x^2 = \mathfrak{p} + x = \bar{x}. \] 
Since $\mathfrak{p}$ is prime, $A / \mathfrak{p}$ is an integral domain. Hence, $\bar{x}^2 = \bar{x} \implies \bar{x} (\bar{x} - 1) = 0 \implies \bar{x} \in \{0, 1\}$. Thus, $A / \mathfrak{p}$ is the field of size 2, making $\mathfrak{p}$ a maximal ideal.
\end{proof}

\begin{proof}[(iii) Proof.]
For $I$ finitely-generated, $I = (x_1,\ldots,x_n)$ for $x_i \in A$. We'll show that $I$ is principal.

We do this by induction, proving that an ideal with $n$-many generators can be generated by a single element. The base case is trivial. For the inductive step,
suppose that for all ideals generated by $n-1$ elements (or fewer), $(a_1,\ldots,a_{n-1}) = (a)$. Consider $(x_1,\ldots,x_n)$. By the inductive step,
$(x_1,\ldots,x_{n-1}) = (a)$, so 
$$(x_1,\ldots,x_n) = (a, x_n).$$

Consider the element $b = a + x_n + ax_n$. Note that 
$$ab = a^2 + ax_n + a^2x_n = a + ax_n + ax_n = a + 2ax_n = a$$
and
$$x_n b = a x_n + x_n^2 + a x_n^2 = a x_n + x_n + a x_n = x_n + 2 a x_n = x_n.$$
Thus, $a, x_n \in (b)$, so $(x_1,\ldots,x_n) = (a, x_n) = (b)$ since, by the construction of $b$, $b \in (a, x_n)$. By induction, any finitely generated
ideal of a Boolean ring $A$ is principal.
\end{proof}

\subsection{Exercise}\label{1.12}

\begin{proof}
If $x$ is a unit, $x^2 = x \implies x^2 x^{-1} = x x^{-1} \implies x = 1$. If $x$ is not a unit, then $x \in \mathfrak{m}$. That means that $1 - x$ is not 
in $\mathfrak{m}$ (otherwise $x + (1-x) = 1 \in \mathfrak{m}$) and consequently that $1 - x$ is a unit. $x^2 = x \implies x (1 - x) = 0$, but since 
$1 - x$ is a unit, we have $x = 0$.
\end{proof}

\subsection{Exercise}\label{1.13}

Here, we construct an algebraic closure of a field $K$.

Let $K$ be a field, and let $\Sigma$ be the set of all irreducible monic polynomials $f$ in one indeterminate with coefficients in $K$. Let 
$A$ be the polynomial ring over $K$ generated by indeterminates $x_f$, one for each $f \in \Sigma$. Let $\mathfrak{a}$ be the ideal of $A$
generated by the polynomials $f(x_f)$ for all $f \in \Sigma$. We now show that $\mathfrak{a} \neq (1)$.

\begin{proof}
Let $R$ be any nonzero commutative ring with unit. Consider $R[x]/(f(x))$, where $f(x)$ is a monic polynomial over $R$. Suppose
$R[x] = (f(x))$. Then, $1 = f(x)r(x)$ for some $r(x) \in R[x]$. However, by Exercise \ref{1.2}(i), $f(x)$ cannot be a unit
since its highest coefficient is $1$, a unit. In particular, $(f(x))$ is a proper ideal of $R[x]$. Thus, for a nonzero ring $R$,
$R[x]/(f(x))$ is also nonzero.

We'll show that $K[x_1,\ldots,x_n] / (f_1(x_1), \ldots, f_n(x_n))$ is nonzero if $\{f_i\}$ are monic polynomials over $K$. We
do this by induction. We proved the base case above, setting $K = R$.

Now, suppose $K[x_1,\ldots,x_{n-1}] / (f_1(x_1), \ldots, f_{n-1}(x_{n-1}))$ is nonzero. Then,
\begin{align*}
K[x_1,\ldots,x_n]/(f_1(x_1), \ldots, f_n(x_n)) &\cong K[x_1,\ldots,x_{n-1}][x_n]/(f_1(x_1), \ldots, f_n(x_n)) \\
&\cong \left(K[x_1,\ldots,x_{n-1}]/(f_1(x_1),\ldots,f_{n-1}(x_{n-1}))\right)[x_n]/(f_n(x_n)) \\
&\cong R[x_n]/(f_n(x_n)),
\end{align*}
where $R = K[x_1,\ldots,x_{n-1}]/(f_1(x_1),\ldots,f_{n-1}(x_{n-1}))$. Since $R$ is nonzero and $f_n(x_n)$ is a monic polynomial over $R$,
we have that $R[x_n]/(f_n(x_n))$ is nonzero by our initial proof. By induction, $K[x_1,\ldots,x_n] / (f_1(x_1), \ldots, f_n(x_n))$ is nonzero.

We turn to the main point. Suppose to the contrary that $\mathfrak{a} = (1)$. Then, there is some finite collection of pairs $(h_i,f_i) \in A \times \Sigma$
such that
$$1 = \sum_{i=1}^n h_i f_i(x_{f_i}).$$
Write $x_i = x_{f_i}$ for convenience. For the indeterminates $x_g$ not equal to some $f_i$, we can map $x_g \mapsto 0$. This is a ring homomorphism of $A$. This
maps the $h_i \mapsto k_i \in K[x_1,\ldots,x_n]$ so that
$$1 = \sum_{i=1}^n k_i f_i(x_i).$$
But then, we have $(f_1(x_1),\ldots,f_n(x_n)) = (1)$, violating $K[x_1,\ldots,x_n]/(f_1(x_1),\ldots,f_n(x_n))$ being nonzero. Thus,
$\mathfrak{a} \neq (1)$. 
\end{proof}

By Corollary 1.4, $A$ has a maximal ideal $\mathfrak{m}$ containing $\mathfrak{a}$ since $\mathfrak{a} \neq (1)$. Now, let $K_1 = A/\mathfrak{m}$,
which is a field since $\mathfrak{m}$ is maximal. Since $\mathfrak{m}$ cannot contain any units and $K \xhookrightarrow{} A$, 
$K$ is a subfield of $K_1$. Further, for every $f \in \Sigma$, $x_f \in K_1$ and satisifies $f(x_f) = 0$.

Since $K_1$ is a field, we can repeat the construction with $K_1$ in place of $K$, obtaining a field $K_2$, and so on. Now, 
let $L = \bigcup_{n=1}^{\infty} K_n$, which is a field since all operations take place in some subfield $K_n$.

Suppose $f \in \Sigma$ has degree $n$. In $K_1$, if $x_f$ maps to the residue class $\alpha_f$ in $A/\mathfrak{m}$, then $f(x) = (x-\alpha_f)g(x)$, where
$g(x) \in K_1[x]$ has degree $n-1$. In the worse case, $g$ is irreducible here. By a simple inductive argument, $f(x)$ factors completely over $K_n \subset K$.
Thus, every $f \in \Sigma$ splits completely in $K$.

By the same logic, let $f \in L[x]$ be a monic polynomial; suppose $f$ has degree $n$. Since $f$ has finitely-many coefficients, all coefficients lie in some
common subfield $K_k$. Then, by construction, $f$ splits completely in $K_{k+n}$, so $L$ is algebraically closed.

Now, let $\overline{K}$ be the set of all elements of $L$ that are algebraic over $K$. Then, $\overline{K}$ is an algebraic closure of $K$.

\subsection{Exercise}\label{1.14}

\begin{proof}
Since $0$ is such an ideal, $\Sigma$ is non-empty. For $I,J \in \Sigma$, we say that $I \le J$ if $I \subset J$, endowing $\Sigma$ with a 
partial ordering. Let $\mathcal{T}$ be a chain of $\Sigma$. Then, for all $I \in \mathcal{T}$, let $T = \bigcup_{I \in \mathcal{T}}I$. $T$ is
an ideal (because the chain is ordered by inclusion) and contains only zero-divisors, so $T \in \Sigma$. Further, $I \subset T$ for all 
$I \in \mathcal{T}$. Thus, every chain has an upper bound in $\Sigma$. By Zorn's lemma, $\Sigma$ has at least one maximal element.

Now that we've proven existence, let $M$ be a maximal element of $\Sigma$. We now show that $M$ is prime. Note that the elements of $A$ that are not
zero divisors form a multiplicatively-closed set $S$. Suppose $x,y \in A - M$. By the maximality of $M$, $M + (x)$ and $M + (y)$ must both contain
elements of $S$. Thus, there are $m_1,m_2 \in M$ and $a_1,a_2 \in A$ so that $m_1 + a_1 x, m_2 + a_2 y \in S$. Since $S$
is multiplicatively-closed,
$$(m_1 + a_1 x)(m_2 + a_2 y) = m_1m_2 + m_1 a_2 y + m_2 a_1 x + a_1 a_2 xy \in S.$$
If $xy \in M$, then all terms would be in $M$ so that the sum would be in $M$. However, $M$ only contains zero divisors and $S$ contains no zero divisors,
so this would be a contradiction. Thus, $xy \not\in M$, proving the primality of $M$.

If $z$ is a zero divisor, then $(z) \in \Sigma$. Thus, $(z)$ is contained in some maximal element $M$ of $\Sigma$, which we've just shown is prime. Thus,
the set of zero divisors in $A$ is a union of prime ideals.
\end{proof}

\subsection{Exercise}\label{1.15}

\begin{proof}[(i) Proof.]
Clearly, $V(\mathfrak{a}) \subset V(E)$ since $E \subset \mathfrak{a}$. However, if $\mathfrak{p}$ contains $E$, then the ideal $\mathfrak{a}$ that $E$ 
generates must also be contained in $\mathfrak{p}$, so $V(E) = V(\mathfrak{a})$.

Since $\mathfrak{a} \subset r(\mathfrak{a})$, we have $V(r(\mathfrak{a})) \subset V(\mathfrak{a})$. Let $\mathfrak{p}$ be any prime containing $\mathfrak{a}$;
consider the ring homomorphism $\varphi: A \to A/\mathfrak{p}$. Then, $\varphi(\mathfrak{a}) = 0$. If $x \in r(\mathfrak{a})$, then $x^n \in \mathfrak{a}$
for some $n > 0$. Under the map, we have $\overline{x}^n = 0$. Since $\mathfrak{p}$ is prime, $A/\mathfrak{p}$ is an integral domain, so $\overline{x} = 0$.
Thus, $x \in \mathfrak{p}$ and $r(\mathfrak{a}) \subset \mathfrak{p}$. Consequently, $V(\mathfrak{a}) = V(r(\mathfrak{a}))$.
\end{proof}

\begin{proof}[(ii) Proof.]
$V(0) = X$ since for any prime ideal $\mathfrak{p} \in X$, $0 \subset \mathfrak{p}$. Likewise, for any prime ideal $\mathfrak{p} \in X$, $\mathfrak{p}$
is a proper ideal of $A$, so $\mathfrak{p}$ does not contain $(1)$. Thus, $V(1) = \emptyset$.
\end{proof}

\begin{proof}[(iii) Proof.]
Let $(E_i)_{i\in I}$ be a family of subsets of $A$. Suppose $\mathfrak{p}$ is a prime ideal containing $\bigcup_{i \in I}E_i$. Then, $\mathfrak{p}$
contains each $E_i$; i.e.,
$$V\left(\bigcup_{i \in I} E_i\right) \subset \bigcap_{i \in I} V(E_i).$$
Conversely, suppose $\mathfrak{p}$ contains each $E_i$; then clearly, it contains the union. So, we have equality.
\end{proof}

\begin{proof}[(iv) Proof.]
If a prime ideal $\mathfrak{p} \in V(\mathfrak{a})\cup V(\mathfrak{b})$, then it is in one of the two sets. Assume without loss of generality that
$\mathfrak{p} \in V(\mathfrak{a})$. Then, $\mathfrak{p}$ contains $\mathfrak{a}$ and thus also contains $\mathfrak{a} \cap \mathfrak{b}$, so
$V(\mathfrak{a}) \cup V(\mathfrak{b}) \subset V(\mathfrak{a}\cap\mathfrak{b})$. Since $\mathfrak{a}\mathfrak{b} \subset \mathfrak{a} \cap \mathfrak{b}$, we have $V(\mathfrak{a} \cap \mathfrak{b}) \subset V(\mathfrak{a}\mathfrak{b})$. Thus,
if we can show that $V(\mathfrak{a}\mathfrak{b}) \subset V(\mathfrak{a}) \cup V(\mathfrak{b})$, we will have shown full equality.

Let $\mathfrak{p}$ be a prime ideal containing $\mathfrak{a}\mathfrak{b}$. Consider the homomorphism $\varphi : A \to A/\mathfrak{p}$. Since
$\varphi$ is surjective, the ideals $\mathfrak{a}$ and $\mathfrak{b}$ map to ideals $\overline{\mathfrak{a}}$ and $\overline{\mathfrak{b}}$ 
of $A/\mathfrak{p}$. Suppose to the contrary that both $\overline{\mathfrak{a}}$ and $\overline{\mathfrak{b}}$ contain nonzero elements, 
say $a$ and $b$. Then, $ab$ is nonzero since $A/\mathfrak{p}$ is an integral domain ($\mathfrak{p}$ is prime). However, 
$\overline{\mathfrak{a}}\cdot\overline{\mathfrak{b}} = \overline{\mathfrak{ab}} = 0$ since $\mathfrak{ab} \subset \mathfrak{p}$. Thus, either $\overline{\mathfrak{a}} = 0$
or $\overline{\mathfrak{b}} = 0$. In other words, $\mathfrak{a} \subset \mathfrak{p}$ or $\mathfrak{b} \subset \mathfrak{p}$. In either case,
$\mathfrak{p} \in V(\mathfrak{a}) \cup V(\mathfrak{b})$, completing the proof.
\end{proof}

\begin{remark}
We've shown that $V(E)$ satisfy the axioms for closed sets in a topological space. The resulting topology is called the Zariski topology. We denote this topological space 
$\Spec(A) := X$, the prime spectrum of $A$.
\end{remark}

\subsection{Exercise}\label{1.16}

\begin{center}
    \includegraphics[width = 4in]{"./images/spec_examples.jpg"}
\end{center}

\subsection{Exercise}\label{1.17}

Before proving properties (i)-(vii), we must show that the sets $X_f$ form a basis of open sets for the Zariski topology.
\begin{proof}
Any open set $U$, by definition, is the complement of a closed set $V$ of $\Spec A$. In Exercise \ref{1.15}, we showed that each closed set of $\Spec A$
can be represented by $V(\mathfrak{a})$ for some ideal $\mathfrak{a}$ (since arbitrary intersections and finite unions have the same form), 
so $U = V(\mathfrak{a})^c$ for some ideal $\mathfrak{a}$. By Exercise \ref{1.15}(iii),
$$V(\mathfrak{a}) = V\left(\bigcup_{f \in \mathfrak{a}}\right) = \bigcap_{f \in \mathfrak{a}}V(f).$$
Thus,
$$U = V(\mathfrak{a})^c = \left(\bigcap_{f \in \mathfrak{a}}V(f)\right)^c
= \bigcup_{f \in \mathfrak{a}}V(f)^c = \bigcup_{f \in \mathfrak{a}}X_f,$$
so $\{X_f\}_{f \in A}$ forms a basis of open sets in $\Spec A$.
\end{proof}

\begin{proof}[(i) Proof.]
Note that $V(f) = V(\{f\})$, the prime ideals containing the element $f \in A$. Then, $X_f = V(\{f\})^c$, where $-^c$ denotes the complement. By Exercise
\ref{1.15}(i), $V(f) = V((f))$ for any $f \in A$. If $f,g \in A$, then using Exercise \ref{1.15}(iv), we have
$$X_f \cap X_g = V((f))^c \cap V((g))^c = (V((f)) \cup V((g)))^c = V((f)(g))^c = V((fg))^c = V(fg)^c = X_{fg},$$
completing the proof.
\end{proof}

\begin{proof}[(ii) Proof.]
If $X_f = \emptyset$, then $V(f) = X$. Thus, $f$ is in the intersection of all prime ideals, meaning $f \in \mathfrak{N}_A$ (i.e., $f$ is nilpotent).
Conversely, if $f$ is nilpotent, it is in the nilradical and thus in every prime ideals, so $X_f = V(f)^c = X^c = \emptyset$.
\end{proof}

\begin{proof}[(iii) Proof.]
If $X_f = X$, then $V(f) = \emptyset$. Thus, $f$ is not in any prime ideal and thus not in any maximal ideal. In particular, $(f) = A$, so $f$ is a unit.
Conversely, if $f$ is a unit, then there is some $g \in A$ such that $fg = 1$. Thus, for any prime ideal $\mathfrak{p}$, under $A \to A/\mathfrak{p}$,
we have $\overline{f}\cdot\overline{g} = 1$. So, $\overline{f} \neq 0$ and $f \not\in \mathfrak{p}$. Hence, $V(f) = \emptyset$ and $X_f = X$.
\end{proof}

\begin{proof}[(iv) Proof.]
If $X_f = X_g$, then $V(f) = V(g)$. By Proposition 1.14,
$$r(f) = \bigcap_{\mathfrak{p} : (f) \subset \mathfrak{p}}\mathfrak{p}
= \bigcap_{\mathfrak{p} \in V(f)}\mathfrak{p}
= \bigcap_{\mathfrak{p} \in V(g)}\mathfrak{p}
= r(g).$$
Conversely, if $r(f) = r(g)$, then $V(f) = V(r(f)) = V(r(g)) = V(g)$, so $X_f = X_g$.
\end{proof}

\begin{proof}[(v) Proof.]
Suppose $\{U_j\}_{j \in J}$ is an open cover of $X$. Since we showed that $\{X_f\}_{f \in A}$ form a basis of open sets, we can write 
$U_j = \bigcup_{f \in I_j}X_f$ for each $j \in J$. Thus, if $I = \bigcup_{j \in J}I_j$, we have an open cover
$$X = \bigcup_{f \in I}X_f = \bigcup_{f \in I}V(f)^c = \left(\bigcap_{f \in I}V(f)\right)^c$$
$$\implies \emptyset = \bigcap_{f \in I}V(f) = V\left(\bigcup_{f \in I}(f)\right),$$
using Exercise \ref{1.15}(iii). Let $\mathfrak{a}$ be the ideal generated by $f \in I$. Then, $V(\mathfrak{a}) = \emptyset$, so 
$V(r(\mathfrak{a})) = \emptyset$ by Exercise \ref{1.15}(i). No prime ideal contains $r(\mathfrak{a})$, so $r(\mathfrak{a}) = (1)$. Thus, $\mathfrak{a} = (1)$.

In particular, there is some finite collection $\{f_s\}_{s \in S}$ of $\{f_i\}_{i \in I}$ such that
$$\sum_{s \in S}a_s f_s = 1$$
for some $a_s \in A$. Thus, there is no prime ideal that divides every $(f_i)$, so $\bigcap_{s \in S}V(f_s) = \emptyset$ and $X = \bigcup_{s \in S}X_{f_s}$.
Since each $f_s$ is contained in some $U_s$, we have
$$X = \bigcup_{s \in S}U_s,$$
where $S$ is finite and $U_s \in \{U_j\}_{j \in J}$. Hence, $X$ is quasi-compact.
\end{proof}

\begin{proof}[(vi) Proof.]
For $f \in A$, let $\{U_j\}_{j \in J}$ be an open covering of $X_f$. Then, for each $j$, 
$U_j = \bigcup_{g \in I_j}X_g$ since $\{X_g\}_{g \in A}$ form a basis of open sets. If $I = \bigcup_{j \in J}I_j$, then
$$X_f \subset \bigcup_{g \in I}X_g.$$
Taking the intersection with $X_f$ and using (i), we have
$$X_f \subset X_f \cap \bigcup_{g \in I}X_g = \bigcup_{g \in I}(X_f \cap X_g) = \bigcup_{g \in I}X_{fg} \subset X_f,$$
so $X_f = \bigcup_{g \in I}X_{fg}$. Taking complements and using results of Exercise \ref{1.15}, we have
$$V(f) = \bigcap_{g \in I} V(fg) = V\left(\bigcup_{g \in I}(fg)\right) = V\left(\sum_{g \in I}(fg)\right).$$
By (iv), $r(f) = r\left(\sum_{g \in I}fg\right)$. Therefore, for some power $n>0$ of $f$, we have
$$f^n = \sum_{s \in S}a_s f g_s$$
for a finite subset $S \subset I$ and $a_s \in A$. From the sum, if there is some prime ideal $\mathfrak{p}$ that does not contain $f$ (and thus not $f^n$),
it does not contain at least one element $a_s f g_s$ on the right-hand side. So, it does not contain $g_s$ for some $s \in S$. In particular,
\begin{align*}
V(f)^c &\subset \bigcup_{s \in S}V(g_s)^c \\
X_f &\subset \bigcup_{s \in S}X_{g_s}.
\end{align*}
Since each $X_{g_s}$ is contained in some $U \in \{U_j\}_{j \in J}$ by the definition of $I$, $X_f$ has a finite subcover and is thus quasi-compact.
\end{proof}

\begin{proof}[(vii) Proof.]
Suppose $U$ is an open set of $X$ and that $U$ is a finite union of sets $X_f$. Now, let $S$ be an open covering of $U$. Then, it
is an open covering of each $X_f$. By (vi), each $X_f$ has a finite subcover. The union of these finite subcovers is finite and covers $U$, so 
$U$ is quasi-compact.

Conversely, suppose $U$ is quasi-compact. By Exercise \ref{1.15}, $U = V(\mathfrak{a})^c$ for some ideal $\mathfrak{a}$. Thus, by Exercise 15(iii),
$$U = V(\mathfrak{a})^c = V\left(\bigcup_{f \in \mathfrak{a}}(f)\right)^c = \left(\bigcap_{f \in \mathfrak{a}}V(f)\right)^c
= \bigcup_{f \in \mathfrak{a}}V(f)^c = \bigcup_{f \in \mathfrak{a}}X_f.$$
Since $U$ is quasi-compact, there is a finite subcover, so $U$ is the union of finitely-many $X_f$.
\end{proof}

\subsection{Exercise}\label{1.18}

\begin{proof}[(i) Proof.]
Let $x \in \Spec A$. If $\mathfrak{p}_x$ is maximal, then $V(\mathfrak{p}_x)$ contains only the prime $\mathfrak{p}_x$. Otherwise, there is a
prime ideal containing $\mathfrak{p}_x$ properly, violating maximality. So in this case, $x$ is a closed point ($\{x\}$ is closed). Conversely,
if $\mathfrak{p}_x$ is not maximal, there is some maximal ideal $\mathfrak{p}_y$ properly containing $\mathfrak{p}_x$. In order for $\{x\}$ to be
closed, there must be some subset $E \subset A$ such that $V(E) = \{x\}$. But in this case, $E \subset \mathfrak{p}_x \subset \mathfrak{p}_y$, so
$\{x,y\}\subset V(E)$. Thus, $x$ is not a closed point when $\mathfrak{p}_x$ is not maximal.
\end{proof}

\begin{proof}[(ii) Proof.]
Any closed set containing $x$ is, by definition, of the form $V(E)$ for some subset $E \subset A$, and by assumption, $x \in V(E)$. Thus, 
$E \subset \mathfrak{p}_x$. By Exercise \ref{1.15}(iii),
$$\overline{\{x\}} = \bigcap_{E \subset \mathfrak{p}_x}V(E) = V\left(\bigcup_{E \subset \mathfrak{p}_x}E\right) = V(\mathfrak{p}_x).$$
\end{proof}

\begin{proof}[(iii) Proof.]
By (ii), $\overline{\{x\}} = V(\mathfrak{p}_x)$, so we need to show that $y \in V(\mathfrak{p}_x) \iff \mathfrak{p}_x \subset \mathfrak{p}_y$. 
But, this holds by definition.
\end{proof}

\begin{proof}[(iv) Proof.]
We aim to show that $\Spec A$ is $T_0$, so for any two elements $x,y \in \Spec A$, we must show that there is an open set that contains only one of
$x$ and $y$. Equivalently, we can produce a closed set that contains exactly one.

Let $x,y \in \Spec A$ be distinct points. If $y \not\in V(\mathfrak{p}_x)$, then we are done since $x \in V(\mathfrak{p}_x)$. So, suppose $y \in V(\mathfrak{p}_x)$. 
By (iii), $\mathfrak{p}_x \subset \mathfrak{p}_y$. Now, consider $V(\mathfrak{p}_y)$. If $x \not\in V(\mathfrak{p}_y)$, then we are done since
$y \in V(\mathfrak{p}_y)$. So, suppose $x \in V(\mathfrak{p}_y)$. By (iii), $\mathfrak{p}_y \subset \mathfrak{p}_x$, so $\mathfrak{p}_x = \mathfrak{p}_y$
(or $x=y$), providing the necessary contradiction.
\end{proof}

\subsection{Exercise}\label{1.19}

\begin{proof}
A topological space is irreducible if $X \neq \emptyset$ and every pair of non-empty open sets in $X$ intersect. Since $\{X_f\}_{f \in A}$ form an open basis of $X$, 
$X$ is irreducible if and only if $X_f \cap X_g \neq \emptyset$ for all $f,g \in A$ satisyfing $X_f \neq \emptyset$ and $X_g \neq \emptyset$.

Suppose that the nilradical $\mathfrak{N}$ is prime, and suppose $f,g \in A$ satisfy $X_f \neq \emptyset$ and $X_g \neq \emptyset$. Then, 
$V(f) \neq X$ and $V(g) \neq X$. Since $\mathfrak{N}$ is the intersect of all primes in $X$, $\mathfrak{N} \not\in V(f) \cup V(g)$ since, by assumption, 
neither $V(f)$ nor $V(g)$ contain all primes. Thus, $\mathfrak{N} \in V(f)^c \cap V(g)^c = X_f \cap X_g$, so $X_f \cap X_g$ is non-empty.

If the nilradical is not prime, then there are $f,g \in A$ such that $fg \in \mathfrak{N}$ yet neither $f$ nor $g$ are nilpotent. By Exercise \ref{1.17}(ii),
$X_f$ and $X_g$ are thus non-empty. However, by Exercise \ref{1.17}(i), $X_f \cap X_g = X_{fg}$. Since $fg \in \mathfrak{N}$ and $\mathfrak{N}$ is the intersection
of all prime ideals, each prime ideal contains $fg$, so $V(fg) = X$ and $X_{fg} = \emptyset$. $X_f \cap X_g = \emptyset$ even though both
$X_f$ and $X_g$ are non-empty, so $X$ is not irreducible.
\end{proof}

\subsection{Exercise}\label{1.20}

\begin{proof}[(i) Proof.]
Let $Y$ be an irreducible subspace of $X$. Because $\overline{Y}$ is closed, any closed set $V$ of $X$ yields a
closed set $V \cap \overline{Y}$ in $X$. Since all closed sets of the subspace $\overline{Y}$ have this form (definition of subspace topology), 
any closed set of $\overline{Y}$ is inherited from a closed set of $X$.

Since $Y$ is irreducible, $Y \neq \emptyset$ (thus $\overline{Y} \neq \emptyset$) and every pair of non-empty open sets in $Y$ intersect. Equivalently, every pair of closed sets $V_1, V_2$ 
not equal to all of $Y$ satisfy $V_1 \cup V_2 \neq Y$.

Let $V_1, V_2$ be any two closed sets of $\overline{Y}$, neither alone containing all of $\overline{Y}$. Then, for some closed sets $A_1,A_2$ of $X$, we have
$V_1 = A_1 \cap \overline{Y}$ and $V_2 = A_2 \cap \overline{Y}$. Thus, $B_1 = A_1 \cap Y$ and $B_2 = A_2 \cap Y$ are closed in $Y$. If, 
for instance, $A_1 \cap Y$ contains all of $Y$, then $A_1 \cap \overline{Y}$ is a closed set containing $Y$, 
so $\overline{Y} \subset A_1 \cap \overline{Y}$. But this is a contradiction since we assumed $A_1 \cap \overline{Y}$ does not contain all of $\overline{Y}$. Thus,
neither $B_1$ nor $B_2$ alone contain all of $Y$. By the irreducibility of $Y$, $B_1 \cup B_2 \neq Y$. 

Now, assume for the sake of contradiction that $V_1 \cup V_2 = \overline{Y}$. Then,
$$\overline{Y} = V_1 \cup V_2 = (A_1 \cap \overline{Y}) \cup (A_2 \cap \overline{Y}) = (A_1 \cup A_2)\cap \overline{Y}.$$
Since $Y \subset \overline{Y}$, we can intersect with $Y$:
$$Y = (A_1 \cup A_2)\cap \overline{Y} \cap Y = (A_1 \cup A_2)\cap Y
= (A_1 \cap Y) \cup (A_2 \cap Y) = B_1 \cup B_2,$$
yielding our contradiction. Thus, $\overline{Y}$ is irreducible when $Y$ is irreducible.
\end{proof}

\begin{proof}[(ii) Proof.]
Let $\Sigma$ be the set of irreducible subspaces of $X$; $\subset$ gives $\Sigma$ a partial ordering. Let $\mathcal{T}$ be a chain, and define 
$U = \bigcup_{Y \in \mathcal{T}} Y$. Let $A_1, A_2$ be two non-empty open sets of $U$. Choose some $x \in A_1$ and $y \in A_2$. 
Then, there is some $Y_1$ that contains $x$ and some $Y_2$ that contains $y$. Since $\mathcal{T}$ is a chain, either $Y_1 \subset Y_2$ or $Y_2 \subset Y_1$. 
Without loss of generality, assume $Y_1 \subset Y_2$. Then, $x \in Y_1 \subset Y_2$, so $x, y \in Y_2$. For convenience, write $Y = Y_2$. 
Note that since $A_1,A_2$ are open in $U$, $A_1 = V_1 \cap U$ and $A_2 = V_2 \cap U$ for open sets $V_1, V_2$ in $X$. But also, $A_1^Y = V_1 \cap Y \subset A_1$ 
and $A_2^Y = V_2 \cap Y \subset A_2$ since $Y \subset U$. Since $x$ is in $A_1$ and $Y$, $x \in A_1^Y$, and since $y$ is in $A_2$ and $Y$, $y \in A_2^Y$, so both sets 
are open and non-empty. By the irreducibility of $Y$, $A_1^Y \cap A_2^Y \neq \emptyset$. Since $A_1^Y \subset A_1$ and $A_2^Y \subset A_2$, $A_1^Y \cap A_2^Y \subset A_1 \cap A_2$,
so $A_1 \cap A_2$ is also non-empty. Thus, $U$ is irreducible and $U \in \Sigma$. In particular, all chains in $\Sigma$ have maximal elements in $\Sigma$. 
By Zorn's lemma, every irreducible subspace of $X$ is contained in a maximal irreducible subspace.
\end{proof}

\begin{proof}[(iii) Proof.]
Let $Y$ be a maximal irreducible subspace of $X$. By (i), $\overline{Y}$ is irreducible, but since $Y$ is maximal and $Y \subset \overline{Y}$, we have
$Y = \overline{Y}$, so $Y$ is closed.

Let $x \in X$; consider the subspace $\{x\}$. If $U_1$ and $U_2$ are any non-empty open sets of $\{x\}$, then $U_1 = U_2 = \{x\}$. Thus, $U_1 \cap U_2 = \{x\} \neq \emptyset$.
In particular, $\{x\}$ is an irreducible subspace. By (ii), it is contained in some maximal irreducible subspace. Thus, the maximal irreducible subspaces cover $X$.

Now, suppose $X$ is Hausdorff. Let $Y$ be an irreducible subspace of $X$, and let $x,y \in Y$ ($Y$ is non-empty by irreducibility). If $x$ and $y$ are distinct, then since $X$ is Hausdorff,
there are open sets $U_x$ and $U_y$ such that $x \in U_x$, $y \not\in U_x$, $y \in U_y$, $x \not\in U_y$, and $U_x \cap U_y = \emptyset$. Thus, $Y \cap U_x$
and $Y \cap U_y$ are non-empty open sets of $Y$ since they contain $x$ and $y$, respectively. However, 
$$(Y \cap U_x) \cap (Y \cap U_y) = Y \cap (U_x \cap U_y) = Y \cap \emptyset = \emptyset.$$
Thus, $Y$ is not irreducible, yielding a contradiction. Thus, $Y$ cannot have more than one point. Conversely, we've already shown that one-point sets are
irreducible subspaces. Thus, the irreducible components of Hausdorff spaces are one-point sets.
\end{proof}

\begin{proof}[(iv) Proof.]
First, we show that $V(\mathfrak{p})$ is irreducible for any $\mathfrak{p} \in X$. In (iii), we showed that one-point sets are irreducible, i.e. $\{\mathfrak{p}\}$ 
is irreducible for every $\mathfrak{p} \in \Spec A$. Since $\overline{\{\mathfrak{p}\}} = V(\mathfrak{p})$ by
Exercise \ref{1.18}(ii), (i) tells us that $\overline{\{\mathfrak{p}\}} = V(\mathfrak{p})$ inherits the irreducibility of $\{\mathfrak{p}\}$.

By (iii), every irreducible component is closed, so each is of the form $V(\mathfrak{a})$ for some ideal $\mathfrak{a}$. If $\mathfrak{a} = (1)$,
then $V(\mathfrak{a}) = \emptyset$, violating irreducibility. Thus, there is some prime ideal $\mathfrak{p}$ such that $\mathfrak{p} \subset \mathfrak{a}$, so
$V(\mathfrak{a}) \subset V(\mathfrak{p})$. Since $V(\mathfrak{p})$ is irreducible, $V(\mathfrak{a}) = V(\mathfrak{p})$ by the maximality of $V(\mathfrak{a})$.
Thus, every irreducible component is of the form $V(\mathfrak{p})$ for some $\mathfrak{p} \in X$.

Now, suppose we have an irreducible component; then, we can write it as $V(\mathfrak{p})$ for some $\mathfrak{p} \in X$. Then, if 
$\mathfrak{q} \in X$ and $\mathfrak{q} \subset \mathfrak{p}$ properly, then $V(\mathfrak{p}) \subset V(\mathfrak{q})$ properly since
$V(\mathfrak{q})$ contains $\mathfrak{q}$ but $V(\mathfrak{p})$ does not contain $\mathfrak{q}$. Since $\mathfrak{q}$ is prime, $V(\mathfrak{q})$ is irreducible,
violating the maximality of $V(\mathfrak{p})$. Thus, there is no such $\mathfrak{q}$, and $\mathfrak{p}$ is a minimal prime ideal of $A$.

Conversely, let $\mathfrak{p}$ be a minimal prime ideal of $A$. Then, $V(\mathfrak{p})$ is contained in some irreducible component, which we've shown 
has the form $V(\mathfrak{q})$ for some $\mathfrak{q} \in X$. Thus, $V(\mathfrak{p}) \subset V(\mathfrak{q})$, so $\mathfrak{q} \subset \mathfrak{p}$.
Since $\mathfrak{p}$ is a minimal prime ideal, $\mathfrak{p} = \mathfrak{q}$, so $V(\mathfrak{p}) = V(\mathfrak{q})$, and $V(\mathfrak{p})$ is an 
irreducible component.
\end{proof}

\subsection{Exercise}\label{1.21}

\begin{proof}[(i) Proof.]
Let $f \in A$. If $\mathfrak{q} \in Y_{\phi(f)}$, then $\mathfrak{q} \not\in V_Y(\phi(f))$, so $(\phi(f)) \not\subset \mathfrak{q}$. Thus,
$\phi(f) \not\in \mathfrak{q}$, so $f \not\in \phi^{-1}(\mathfrak{q})$. Hence, $\phi^{-1}(\mathfrak{q}) \not\in V_X(f)$, so $\phi^{-1}(\mathfrak{q}) \in X_f$ 
and $\mathfrak{q} \in \phi^{*-1}(X_f)$.

Conversely, if $\mathfrak{q} \not\in Y_{\phi(f)}$, then $\mathfrak{q} \in V_Y(\phi(f))$, so $\phi(f) \in \mathfrak{q}$. Thus, $f \in \phi^{-1}(\mathfrak{q})$, so
$\phi^{-1}(\mathfrak{q}) \in V_X(f)$. Hence, $\phi^{-1}(\mathfrak{q}) \not\in X_f$. Then, by definition, $\mathfrak{q} \not\in \phi^{*-1}(X_f)$. Thus, $\phi^{*-1}(X_f) = Y_{\phi(f)}$ for any $f \in A$.

Let $U$ be an open set of $X$, so $U = \bigcup_{f \in I}X_f$. Then,
$$\phi^{*-1}(U) = \phi^{*-1}\left(\bigcup_{f \in I}X_f\right) = \bigcup_{f \in I}\phi^{*-1}(X_f) = \bigcup_{f \in I}Y_{\phi(f)},$$
which is an open set of $Y$. Thus, $\phi^* : Y \to X$ is continuous.
\end{proof}

\begin{proof}[(ii) Proof.]
If $\mathfrak{q} \in \phi^{*-1}(V_X(\mathfrak{a}))$, then $\phi^*(\mathfrak{q}) = \phi^{-1}(\mathfrak{q}) \in V_X(\mathfrak{a})$. Hence,
$\mathfrak{a} \subset \phi^{-1}(\mathfrak{q})$, so $\phi(\mathfrak{a}) \subset \mathfrak{q}$, so 
$\mathfrak{q} \in V_Y(\phi(\mathfrak{a})) = V_Y(\mathfrak{a}^e)$.

Conversely, if $\mathfrak{q} \in V_Y(\mathfrak{a}^e)$, then $\phi(\mathfrak{a}) \subset \mathfrak{a}^e \subset \mathfrak{q}$, so 
$\mathfrak{a} \subset \phi^{-1}(\mathfrak{q})$, hence $\phi^{-1}(\mathfrak{q}) \in V_X(\mathfrak{a})$. Thus, $\mathfrak{q} \in \phi^{*-1}(V_X(\mathfrak{a}))$. 
Combining both directions, we have $\phi^{*-1}(V_X(\mathfrak{a})) = V_Y(\mathfrak{a}^e)$.
\end{proof}

\begin{proof}[(iii) Proof.]
Note that $\phi^*(V(\mathfrak{b})) = \{\mathfrak{q}^c : \mathfrak{b}\subset\mathfrak{q}\}$. So, $\mathfrak{a}$ must contain all of these 
$\mathfrak{q}^c$. Since all closed sets containing $\phi^*(V(\mathfrak{b}))$ must contain each $\mathfrak{q}^c$, all must contain 
$\bigcap_{\mathfrak{b} \subset \mathfrak{q}}\mathfrak{q}^c$. Since this is itself an ideal, it gives the closure:
$$\overline{\phi^*(V(\mathfrak{b}))} = V\left(\bigcap_{\mathfrak{b} \subset \mathfrak{q}}\mathfrak{q}^c\right).$$
Note that 
$$\bigcap_{\mathfrak{b} \subset \mathfrak{q}}\mathfrak{q}^c = \bigcap_{\mathfrak{b} \subset \mathfrak{q}}\phi^{-1}(\mathfrak{q})
= \phi^{-1}\left(\bigcap_{\mathfrak{b} \subset \mathfrak{q}}\mathfrak{q}\right) = \phi^{-1}(r(\mathfrak{b})) = r(\mathfrak{b})^c = r(\mathfrak{b}^c).$$
Thus,
$$\overline{\phi^*(V(\mathfrak{b}))} = V\left(\bigcap_{\mathfrak{b} \subset \mathfrak{q}}\mathfrak{q}^c\right) = V(r(\mathfrak{b}^c)) = V(\mathfrak{b}^c).$$
\end{proof}

\begin{proof}[(iv) Proof.]
Since $\phi$ is surjective, we have a ring isomorphism $A/(\ker \phi) \to B$. Thus, we have an inverse homomorphism 
$\psi : B \to A/(\ker \phi)$. From these maps, we obtain $\psi^* : \Spec (A/(\ker \phi)) \to \Spec (B)$ and 
$\phi^* : \Spec(B) \to \Spec(A/(\ker \phi))$, which are inverse maps describing a bijection; by (i), both are continuous, so they describe a homeomorphism.

Now, we have to show that the map $\Spec(A / (\ker \phi)) \to V(\ker \phi)$, viewed as a closed subspace of $\Spec(A)$, is a homeomorphism.

But this follows easily from Proposition 1.1, considering the map $A \to A/(\ker \phi)$. More concretely, 
there is a one-to-one, order-preserving, primality-preserving correspondence between ideals of $A$ containing $\ker \phi$
and ideals of $A/(\ker \phi)$. Thus, the map induced by projection $\Spec(A/\ker \phi) \to V(\ker\phi)$ is a bijection. 
Because the map preserves the containment relation, it is a closed map ($V(\mathfrak{a}) \to V(\mathfrak{a}^c)$). Since it is a bicontinuous bijection,
it is a homeomorphism.
\end{proof}

\begin{remark}
This shows that $\Spec(A)$ and $\Spec(A/\mathfrak{N}_A)$ are naturally homeomorphic.
\end{remark}

\begin{proof}[(v) Proof.]
Suppose $\ker \phi \subset \mathfrak{N}$. Since $Y = V(0)$, by (iii), we have 
$$\overline{\phi^*(Y)} = \overline{\phi^*(V(0))} = V(\phi^{-1}(0)).$$ Since $\ker \phi$ is in the nilradical, which is contained by every prime ideal,
$V(\phi^{-1}(0)) = X$, so $\phi^*(Y)$ is dense in $X$. 

Conversely, if $\phi^*(Y)$ is dense in $X$, then we have, by (iii),
$$X = \overline{\phi^*(Y)} = \overline{\phi^*(V(0))} = V(\phi^{-1}(0)),$$
so $\phi^{-1}(0) \subset \mathfrak{p}$ for all prime ideals, so $\ker \phi \subset \mathfrak{N}$.
\end{proof}

\begin{proof}[(vi) Proof.]
We have two ring homomorphisms $\phi : A \to B$ and $\psi : B \to C$. If $\mathfrak{r} \in \Spec(C)$, then we have a map $(\psi \circ \phi)^* : \Spec(C) \to \Spec(A)$. This takes
$\mathfrak{r} \in \Spec(C)$ to $(\psi \circ \phi)^{-1}(\mathfrak{r}) \in \Spec(A)$. But, $(\psi \circ \phi)^{-1} = \phi^{-1} \circ \psi^{-1}$. In particular,
$(\psi \circ \phi)^*$ decomposes into maps 
$$\Spec(C) \xlongrightarrow{\psi^*} \Spec(B) \xlongrightarrow{\phi^*} \Spec(A),$$ 
so $(\psi \circ \phi)^* = \phi^* \circ \psi^*$.
\end{proof}

\begin{proof}[(vii) Proof.]
$A$ is an integral domain with one non-zero prime ideal $\mathfrak{p}$, $K$ is the field of fractions of $A$, and $B = (A/\mathfrak{p}) \times K$. We
define $\phi: A \to B$ by $\phi(x) = (\overline{x}, x)$, where $\overline{x}$ is the image of $x$ in $A/\mathfrak{p}$.

Since $\phi$ is a ring homomorphism, we have $\phi^*: \Spec B \to \Spec A$. $A$ has only one non-zero prime ideal $\mathfrak{p}$, so $\Spec A = \{0,\mathfrak{p}\}$.
For $B$, the only prime ideal of $K$ is $0$, and the only prime ideal of $A / \mathfrak{p}$ is $0$. Thus, $\Spec B = \{0 \times K, (A / \mathfrak{p}) \times 0\}$. Note
that $0$ is not prime in $B$ since $0 = (0 \times K)((A / \mathfrak{p}) \times 0)$. We have $\phi^{-1}(0 \times K) = \mathfrak{p}$ and $\phi^{-1}((A / \mathfrak{p}) \times 0) = 0$.
Thus, $\phi^*$ is bijective.

However, we will show that it is not a homeomorphism. We only need to look at the image $\phi^*(\{(A / \mathfrak{p}) \times 0\}) = \{0\}$. Note that 
$\{(A / \mathfrak{p}) \times 0\} = V((A / \mathfrak{p}) \times 0)$, so this set is closed. However, $\{0\}$ is not closed since for any $\mathfrak{a}$, 
$0 \in V(\mathfrak{a})$ implies $\mathfrak{a} \subset 0 \subset \mathfrak{p}$, so $\mathfrak{p} \in \overline{\{0\}}$, i.e. $\phi^*$ maps a closed set to a non-closed
set. Therefore, $\phi^*$ is not a homeomorphism.
\end{proof}

\subsection{Exercise}\label{1.22}

First, we show that if $A = \prod_{i=1}^n$, then $\bigsqcup_{i=1}^n \Spec(A_j) \cong \Spec A$ as topological spaces.

\begin{proof}
Let $A = \prod_{i=1}^n$. We claim that all primes take the form 
$$A_1 \times \cdots \times A_{j-1} \times \mathfrak{p}_j \times A_{j+1} \times \cdots \times A_n,$$
where $\mathfrak{p}_j \in \Spec A_j$. If $\pi_j : A \to A_j$ is the projection map, this is equivalent to saying all primes are of the form
$\pi_j^{-1}(\mathfrak{p}_j)$ for some $j$ and $\mathfrak{p}_j \in \Spec(A_j)$. First of all, it's clear that ideals of this form 
are prime since for $a,b \in A$ with $\pi_j(ab) \in \mathfrak{p}$, we have $\pi_j(a)\pi_j(b) \in \mathfrak{p}$, so either $\pi_j(a) \in \mathfrak{p}$
or $\pi_j(b) \in \mathfrak{p}$. Thus, either $a \in \pi_j^{-1}(\mathfrak{p})$ or $b \in \pi_j^{-1}(\mathfrak{p})$, so $\pi_j^{-1}(\mathfrak{p}) \in \Spec(A)$.
In particular, we've shown $$\bigsqcup_{i=1}^n \Spec(A_j) \xhookrightarrow{} \Spec A.$$

Conversely, suppose $\mathfrak{p}$ is a prime of $A$. Suppose $\pi_i(\mathfrak{p}) \neq A_i$ and $\pi_j(\mathfrak{p}) \neq A_j$ for distinct $i,j$. Then,
let the components of $a \in A$ be all zeros except for the $i$th slot, where we choose $a_i \in A_i - \pi_i(\mathfrak{p})$. Similarly, we define $b \in A$ by
letting the components of $b$ be all zeros except for the $j$th slot, where we choose $b_j \in A_j - \pi_j(\mathfrak{p})$. Then, neither $a$ nor $b$ is in $\mathfrak{p}$,
yet $ab = 0 \in \mathfrak{p}$, violating the primality of $\mathfrak{p}$. Thus, all projection maps of $\mathfrak{p}$ must be surjective except for one. So,
$$\mathfrak{p} = A_1 \times \cdots \times A_{j-1} \times \pi_j(\mathfrak{p}) \times A_{j+1} \times \cdots \times A_n.$$
Let $a_j,b_j \in A_j$. Lift $a_j,b_j$ to $a,b \in A$ by letting all other components be zero. If $ab \in \mathfrak{p}$, then $a_jb_j \in \pi_j(\mathfrak{p})$. Since
$\mathfrak{p}$ is prime, either $a$ or $b$ is in $\mathfrak{p}$. Without loss of generality, suppose $a \in \mathfrak{p}$. Then, $\pi_j(a) = a_j \in \pi_j(\mathfrak{p})$. Thus,
$\pi_j(\mathfrak{p})$ is prime. Hence, we have a bijection
$$\bigsqcup_{i=1}^n \Spec(A_j) \cong \Spec A.$$

To be more concrete, consider the projection $\pi_j : A \to A_j$. By Exercise \ref{1.21}(iv), $\pi_j^*$ is a homeomorphism of $\Spec(A_j)$ onto the closed subset
$V(\ker \pi_j)$ of $\Spec A$. Explicitly, the kernel is $A_1 \times \cdots \times A_{j-1} \times 0 \times A_{j+1} \times \cdots \times A_n$, which matches
our constructed maps. We identified $V(A_1 \times \cdots \times A_{j-1} \times 0 \times A_{j+1} \times \cdots \times A_n)$ directly with $\Spec A_j$, which are homeomorphic. 
To be slightly pedantic, note that the closed subsets just mentioned are disjoint for distinct $j$ by our earlier proof that a prime cannot have two non-surjective projections.
\end{proof}

\begin{proof}
By the previous proof, (ii) implies (i) since $A \cong A_1 \times A_2$ implies $\Spec(A) \cong \Spec(A_1) \sqcup \Spec(A_2)$, so $\Spec(A)$ is disconnected. 

Now, we show (i) implies (iii). Suppose $\Spec(A)$ is disconnected. By Exercise \ref{1.21}(iv), $\Spec(A)$ is naturally homeomorphic to $\Spec(A/\mathfrak{N})$, 
so consider $A' = A/\mathfrak{N}$, which has no nilradical and is disconnected (by the homeomorphism). We can write $\Spec(A')$ as the disjoint union 
of two closed sets $V(\mathfrak{a})$ and $V(\mathfrak{b})$  for ideals $\mathfrak{a},\mathfrak{b} \subset A'$. By Exercise \ref{1.15}(iv),
$$\Spec(A') = V(\mathfrak{a}) \cup V(\mathfrak{b}) = V(\mathfrak{a} \cap \mathfrak{b}),$$
so $\mathfrak{a} \cap \mathfrak{b} \subset \mathfrak{N}_{A'} = 0$. Further, since $V(\mathfrak{a}) \cap V(\mathfrak{b}) = \emptyset$, $\mathfrak{a} + \mathfrak{b} = A'$.
Consider the ring homomorphism $\psi: A' \to A'/\mathfrak{a} \times A'/\mathfrak{b}$. If
$\psi(x) = (0,0)$, then $x \in \mathfrak{a} \cap \mathfrak{b} \subset \mathfrak{N}_{A'} = 0$, so $\psi$ is injective. By Proposition 1.10(ii), since
$\mathfrak{a} + \mathfrak{b} = (1)$, $\psi$ is surjective, so we have a ring isomorphism. Thus, 
$$A / \mathfrak{N} \cong A_1 \times A_2$$ for nonzero rings $A_1$ and $A_2$. Now, consider the idempotent $\overline{e} = (1,0)$. Then, there is an $e \in A$ such 
that $e^2 - e \in \mathfrak{N}$. If $e^2 - e = 0$, then we have our idempotent; otherwise, there is a minimal $n > 0$ such that
$(e^2 - e)^n = 0$. We can use the identity in $\Z[X]$ (the binomial theorem), which expands $1 = (X + (1-X))^{2n-1}$ as
$$1 = \sum_{k=0}^{2n-1}\binom{2n-1}{k}e^k (1-e)^{2n-1-k}.$$
Let $f = \sum_{k=0}^{n-1}\binom{2n-1}{k}e^k(1-e)^{2n-1-k}$. Then, we have
$$f = f^2 + f\sum_{k=n}^{2n-1}\binom{2n-1}{k}e^k(1-e)^{2n-1-k},$$
but every term of $f$ contains $(1-e)$ of degree at least $n$, and all terms of the remainder contain $e$ of degree at least $n$. Thus, the latter sum is annihilated,
and we have $f = f^2$. Note that since $e(1-e) \in \mathfrak{N}$, $f \equiv (1-e)^{2n-1} \mod \mathfrak{N}$. Projecting down to $A_1 \times A_2$, we have
$\overline{f} = (1 - (1,0))^{2n-1} = (0,1)^{2n-1} = (0,1)$, which is not the image of either $0$ or $1$. Thus, as desired, $f$ is a non-trival idempotent.

Now, we show (iii) implies (ii). Suppose $e$ is an idempotent of $A$ not equal to $0$ or $1$. Note that
$$(1-e)^2 = 1 - 2e + e^2 = 1 - 2e + e = 1 - e,$$
so $1-e$ is also an idempotent. Further, note that $e(1-e) = e - e^2 = e - e = 0$, so $e$ and $1-e$ are orthogonal idempotents. For any $x \in A$, we have
$x = xe + x(1-e)$. Then, we have the decomposition $\psi : A \cong Ae \times A(1-e)$ by multiplying by $e$ and $1-e$, respectively, i.e. $\psi(x) = (xe,x(1-e))$. 
Note that the right-hand side is a ring since $e$ acts as the identity in $Ae$ and $(1-e)$ acts as the identity in $A(1-e)$. Further, it's easily checked that the 
map is a ring homomorphism by using the idempotent relation and the orthogonality of the two idempotents. The only thing left to check is that the map is bijective.
If $\psi(x) = (0,0)$, then $xe=0$ and $x(1-e) = 0$, so their sum $xe + x(1-e) = x = 0$, so injectivity holds. Conversely, if we have $(ae,b(1-e))$ in the direct product, 
then $ae + b(1-e) \in A$ maps to this element, so surjectivity holds. Thus, this is an isomorphism of rings. Finally, (iii) implies (ii), and now we see that 
all three properties are equivalent.
\end{proof}

\begin{remark}
By Exercise \ref{1.12}, no local ring has an idempotent distinct from $0$ and $1$. We've shown this cannot happen if $\Spec(A)$ is disconnected, so the spectrum of
a local ring is always connected.
\end{remark}

\subsection{Exercise}\label{1.23}

\begin{proof}[(i) Proof.]
For any $f \in A$, consider $X_f$. Clearly, $X_f$ is closed and open if and only if $V(f)$ is open and closed. Since $V(f)$ is closed by definition, we just have to show that $V(f)$
is an open set, i.e. there is some $g \in A$ such that $V(f) = V(g)^c$. Since $A$ is a boolean ring, $f = f^2$ and $(1-f)^2 = 1-f$. Note that for any $x \in A$,
we can write $x = xf + x(1-f)$. Consider the map  $\psi : A \to Af \times A(1-f)$. We will show that this is a ring isomorphism. It's a homomorphism since $Af$ and $A(1-f)$ are rings
with multiplicative identity $f$ and $1-f$; also, since $f(1 - f) = f - f^2 = 0$, multiplication is respected, i.e. 
$(af + b(1-f))(cf + d(1-f)) = (ac)f + (bd)(1-f)$. Now, we must show it is bijective. If $\psi(x) = (0,0)$, then $xf = 0$ and $x(1-f) = 0$, so $0 = xf + x(1-f) = x$, hence $\psi$ is 
injective. Conversely, if we have $(af, b(1-f)) \in Af \times A(1-f)$, then $\psi(af + b(1-f)) = (af, b(1-f))$. Thus, $\psi$ is surjective and hence bijective.

Since $\psi$ realizes $A$ as a direct product of rings, by Exercise \ref{1.22}, we know that 
$$\Spec(A) \cong \Spec(Af) \sqcup \Spec(A(1-f)).$$
Now, considering the projections $\pi_f : A \to Af$ and $\pi_{1-f} A \to A(1-f)$, we have by Exercise \ref{1.21}(iv) that $\Spec(Af) \cong V(f)$ and $\Spec(A(1-f)) \cong V(1-f)$, where $\cong$
denotes homeomorphism. Thus, $\Spec(A)$ is a disjoint union of the closed sets $V(f)$ and $V(1-f)$, so as desired, $V(f) = V(1-f)^c$.
\end{proof}

\begin{proof}[(ii) Proof.]
Let $f_1,\ldots,f_n \in A$. We want to show that $X_{f_1} \cup \cdots \cup X_{f_n} = X_f$ for some $f \in A$. Taking complements, this is equivalent to the statement
$$V(f_1) \cap \cdots \cap V(f_n) = V(f).$$
Note that by Exercise \ref{1.15}(iii), we have
$$V(f_1) \cap \cdots \cap V(f_n) = V\left(\bigcup_{i=1}^n (f_i)\right) = V((f_1,\ldots,f_n)).$$
However, by Exercise \ref{1.11}(i), every finitely generated ideal of $A$ is principal, so $(f_1,\ldots,f_n) = (f)$ for some $f \in A$, as desired.
\end{proof}

\begin{proof}[(iii) Proof.]
Let $Y \subset X$ be a closed and open set. Since $Y$ is closed and $X$ is quasi-compact by Exercise \ref{1.17}(v), we can easily show that $Y$ is also quasi-compact. Let $\{U\}_U$ be an open cover of 
$Y$. Then, $Y^c$ is open, so $\{Y^c\}\cup\{U\}_U$ is an open cover of $X$. Since $X$ is quasi-compact, there is a finite subcover $\{V\}_V$, which obviously covers $Y$ as well. Removing $Y^c$ if necessary, which
does not cover any points of $Y$, we have a finite subcover of $Y$ that uses only open sets in $\{U\}_U$. Thus, $Y$ is quasi-compact.

Since $Y$ is also open, it is a union of basic open sets, and since it is quasi-compact, it is a union of finitely-many basic open sets. In other words, $Y = \bigcup_{i=1}^n X_{f_i}$. By (ii), there is some $f \in A$
such that $\bigcup_{i=1}^n X_{f_i} = X_f$, so $Y = X_f$, as desired.
\end{proof}

\begin{proof}[(iv) Proof.]
Let $\mathfrak{p},\mathfrak{q}$ be distinct primes in $\Spec A$. Since they are distinct, suppose without loss of generality that $\mathfrak{p}$ contains some $f \in A$ that is not in $\mathfrak{q}$; otherwise, swap $\mathfrak{q}$
and $\mathfrak{p}$. Then, $(f) \subset \mathfrak{p}$, so $\mathfrak{p} \not\in X_f$. But, since $\mathfrak{q} \not\in V(f)$, we have $\mathfrak{q} \in X_f$. Since $X_f$ is both open and closed by (i), we have two disjoint open 
sets $X_f$ and $X_f^c$ that contain $\mathfrak{q}$ and $\mathfrak{p}$, respectively. Thus, $\Spec A$ is Hausdorff. Since $X$ is quasi-compact by Exercise \ref{1.17}(v), $X$ is a compact Hausdorff space. 
\end{proof}

\subsection{Exercise}\label{1.24}

\begin{proof}
We start by showing $A(L)$ is a ring. We'll first prove the useful fact:
$$(a \wedge b)' = a' \vee b'.$$
Since the complement is unique, we just have to verify that it is the complement of $a \wedge b$. We have
\begin{align*}
(a \wedge b) \wedge (a' \vee b') &= a \wedge b \wedge (a' \vee b')\\
&= a \wedge ((b \wedge a') \vee (b \wedge b'))\\
&= a \wedge ((b \wedge a') \vee 0)\\
&= a \wedge (b \wedge a')\\
&= (a \wedge a') \wedge b \\
&= 0 \wedge b \\
&= 0.
\end{align*}
Then, we have
\begin{align*}
(a \wedge b) \vee (a' \vee b') &= (a \wedge b) \vee a' \vee b'\\
&= ((a \vee a') \wedge (b \vee a')) \vee b'\\
&= (1 \wedge (b \vee a')) \vee b'\\
&= (b \vee a') \vee b'\\
&= a' \vee (b \vee b')\\
&= a' \vee 1\\
&= 1,
\end{align*}
so indeed, $(a \wedge b)' = a' \vee b'$. In fact, we can permute $a \maspto a'$ and $b \mapsto b'$ and take the complement of both sides to obtain the complementary property:
$$a' \wedge b' = (a \vee b)'.$$

Now, we can begin proving that $A(L)$ is a ring, starting with the additive properties. Note that because the complements are unique, we have $0' = 1$ since $0 \wedge 1 = 0$ and $0 \vee 1 = 1$. Also note that
since $a$ and $a'$ are symmetric in the definition of $a'$, we have $(a')' = a$.

By the commutativity of $\vee$ and $\wedge$, $+$ is also commutative. We'll now show associativity.
\begin{align*}
(a + b) + c &= ((a \wedge b') \vee (a' \wedge b)) + c \\
&= (((a \wedge b') \vee (a' \wedge b)) \wedge c') \vee (((a \wedge b') \vee (a' \wedge b))' \wedge c) \\
&= ((a \wedge b' \wedge c') \vee (a' \wedge b \wedge c')) \vee (((a \wedge b') \vee (a' \wedge b))' \wedge c) \\
&= (a \wedge b' \wedge c') \vee (a' \wedge b \wedge c') \vee (((a \wedge b')' \wedge (a' \wedge b)') \wedge c) \\
&= (a \wedge b' \wedge c') \vee (a' \wedge b \wedge c') \vee ((a' \vee b) \wedge (a \vee b') \wedge c) \\
&= (a \wedge b' \wedge c') \vee (a' \wedge b \wedge c') \vee (a' \wedge a \wedge c) \vee (a' \wedge b' \wedge c) \vee (b \wedge a \wedge c) \vee (b \wedge b' \wedge c) \\
&= (a \wedge b' \wedge c') \vee (a' \wedge b \wedge c') \vee (0 \wedge c) \vee (a' \wedge b' \wedge c) \vee (b \wedge a \wedge c) \vee (0 \wedge c) \\
&= (a \wedge b' \wedge c') \vee (a' \wedge b \wedge c') \vee 0 \vee (a' \wedge b' \wedge c) \vee (b \wedge a \wedge c) \vee 0 \\
&= (a \wedge b' \wedge c') \vee (a' \wedge b \wedge c') \vee (a' \wedge b' \wedge c) \vee (a \wedge b \wedge c).
\end{align*}
Note that the final line is completely symmetric under permutations of $a$, $b$, and $c$. Since we already have commutativity, this is equal to $c + (b + a)$. Since permuting $a$ and $c$ does not
change the value of the final line (by aforementioned symmetry), we have $(a + b) + c = a + (b + c)$, so addition is associative.

For any $a \in A(L)$, we have
$$a + 0 = (a \wedge 0') \vee (a' \wedge 0) = (a \wedge 1) \vee (a' \wedge 0) = a \vee 0 = a,$$
so $0$ is the additive identity. Further, for any $a \in A(L)$, we have
$$a + a' = (a \wedge a') \vee (a \wedge a') = a \wedge a' = 0,$$
by definition, so each $a$ has an additive inverse. Thus, $(A(L), +)$ forms an abelian group.

Moving on, note that multiplication is commutative because $\wedge$ is commutative. Then, let $a,b \in A(L)$. We have $a \cdot 1 = a \wedge 1 = a$ since $1$ is a global max. Thus, $A(L)$ has commutative
multiplication with unit. Associativity follows easily since
$$a\cdot(b\cdot c) = a \wedge(b \wedge c) = (a \wedge b) \wedge c = (a \cdot b) \cdot c.$$
Lastly, we must show that multiplication distributes over addition. We have
\begin{align*}
ab + ac &= (a \wedge b) + (a \wedge c) \\
&= ((a \wedge b) \wedge (a \wedge c)') \vee ((a \wedge b)' \wedge (a \wedge c)) \\ 
&= (a \wedge b \wedge (a' \vee c')) \vee ((a' \vee b') \wedge a \wedge c) \\ 
&= (a \wedge b \wedge a') \vee (a \wedge b \wedge c') \vee (a'\wedge a \wedge c) \vee (b'\wedge a \wedge c) \\ 
&= (0 \wedge b) \vee (a \wedge b \wedge c') \vee (0 \wedge c) \vee (b'\wedge a \wedge c) \\ 
&= 0 \vee (a \wedge b \wedge c') \vee 0 \vee (b'\wedge a \wedge c) \\ 
&= (a \wedge b \wedge c') \vee (a \wedge b' \wedge c) \\ 
&= a \wedge ((b \wedge c') \vee (b' \wedge c)) \\ 
&= a \wedge (b + c) \\
&= a(b+c), 
\end{align*}
so the distributive property holds, and $A(L)$ is a ring. It is now easy to see that it is a Boolean ring since for any $a \in A(L)$, $a^2 = a \wedge a = a$.

Conversely, suppose $A$ is a Boolean ring. We define an ordering $a \le b$ when $a = ab$. Thus, $0$ is clearly a minimal element and $1$ is clearly a maximal element,
satisfying property (i) of a Boolean lattice. Define $a \wedge b = ab$ and $a \vee b = a + b + ab$. We now verify the two distributive properties:
\begin{align*}
a \wedge (b \vee c) &= a(b + c + bc) \\
&= ab + ac + abc \\
&= ab + ac + a^2bc \\
&= ab + ac + (ab)(ac) \\
&= (ab) \vee (ac) \\
&= (a \wedge b) \vee (a \wedge c).
\end{align*}
Also,
\begin{align*}
(a \vee b) \wedge (a \vee c) &= (a + b + ab)(a + c + ac) \\
&= a^2 + ac + a^2c + ab + bc + abc + a^2b + abc + a^2bc \\
&= a + ac + ac + ab + bc + abc + ab + abc + abc \\
&= a + 2ac + 2ab + bc + abc + 2abc \\
&= a + bc + abc \\
&= a \vee (bc) \\
&= a \vee (b \wedge c),
\end{align*}
so property (ii) is satisfied.

Lastly, we must show that there is a unique complement $a'$ for each $a \in A$. We claim that this is satisfied by $1 - a$. We can prove existence and uniqueness at once by
applying the definitions. We have
$$0 = a \wedge a' = a a' \quad\text{ and }\quad 1 = a \vee a' = a + a' + a a'.$$
Since $a a' = 0$, we have that $1 = a + a' + 0$, so $a' = 1 - a$. This proves (iii). Thus, $A$ defined by our $\wedge$ and $\vee$ is a Boolean lattice.
\end{proof}

\subsection{Exercise}\label{1.25}

\begin{proof}
Here, we prove Stone's theorem. Let $L$ be a Boolean lattice. By Exercise \ref{1.24}, $L$ can be made into a Boolean ring, $A(L)$. We define $\mathcal{T}$ to be the open-and-closed subsets of $\Spec(A(L))$. 
By Exercise \ref{1.23}, the elements of $\mathcal{T}$ are exactly $X_f$, where $f \in A(L)$. By Exercise \ref{1.23}, $\Spec(A(L))$ is a compact Hausdorff space, so $\mathcal{T}$ is the set of open-and-closed subsets
of a compact Hausdorff space. Thus, we aim to show that $L \to \mathcal{T}$ is an isomorphism of lattices.

We show that $\mathcal{T}$ is a Boolean lattice. We define $U \le V$ when $U \subset V$. (i) $X_1$ is the maximal element; $X_0$ is the minimal element. (ii) We define $\wedge$ and $\vee$ by $\cap$ 
and $\cup$. These are closed operations (two open-and-closed sets yield an open-and-closed set), and they distribute over each other. (iii) We define the complement of $X_f$ by $X_f^c = X_{f'}$; 
there is such an $f'$ since $X_f^c$ is open-and-closed. $X_f \wedge X_{f'} = X_f \cap X_{f'} = \emptyset = X_0$ and $X_f \vee X_{f'} = X_f \cup X_{f'} = X = X_1$. Thus, $\mathcal{T}$ is a lattice.

The map in view factors as $L \to A(L) \to \mathcal{T}$, where the first map is clearly a bijection. Suppose $f,g \in A(L)$. Then, if $F,G \in L$ yield $f,g \in A(L)$, then $fg = F \wedge G$ by
our construction of the Boolean ring (Exercise \ref{1.24}). Thus, $F \wedge G \mapsto fg \mapsto X_{fg}$. However, $F \mapsto f \mapsto X_f$ and $G \mapsto g \mapsto X_g$. However, $X_f \wedge X_g = X_f \cap X_g = X_{fg}$
by Exercise \ref{1.17}, so the map preserves $\wedge$. If $f, f'$ are complements in $L$, then in $A(L)$, they satisfy 
$$f + f' = (f \wedge f) \vee (f' \wedge f') = f \vee f' = 1,$$
so $f' = 1-f$. However, $(f) + (1-f) = (1)$ and $(f) \cap (1-f)$ contains no prime ideals. Thus, $X_f^c = X_{1-f}$, so the map preserves complements. Because the map preserves complements and $\wedge$,
de Morgan's law here shows that it also preserves $\vee$. In particular, $L \to \mathcal{T}$ is a homomorphism of lattices.

Lastly, we show bijectivity. Note that $L \to A(L)$ is a bijection since the map preserves the underlying set. Every element of $\mathcal{T}$ has the form $X_f$ for some $f \in A$, and the map is $f \mapsto X_f$,
so clearly the map is surjective. To show injectivity, suppose $f,g \in X_f$ and $X_f = X_g$. Then, for every prime ideal $\mathfrak{p}$ of $A(L)$, $f \in \mathfrak{p} \iff g \in \mathfrak{p}$. By Exercise \ref{1.11},
$A/\mathfrak{p}$ is the field of two elements, so for all $\mathfrak{p}$, $f \equiv g \mod \mathfrak{p}$. In particular, $f-g \in \mathfrak{p}$ for all $\mathfrak{p}$. Hence, $f-g \in \mathfrak{N}$. Since
$a^2 = a$ for all $a \in A(L)$ (Boolean ring), $(f-g)^2 = f-g$. However, $f-g$ is nilpotent, so $(f-g)^n = 0$ for some $n > 0$. By a simple inductive argument using $(f-g)^2 = f-g$, we have $f-g = (f-g)^n = 0$,
so $f = g$. Thus, the map is injective and therefore bijective. Hence, $L \to \mathcal{T}$ is a lattice isomorphism.
\end{proof}

\subsection{Exercise}\label{1.26}

Let $A$ be a ring. We define $\Max(A)$ to be the subspace of $\Spec(A)$ containing the maximal ideals, and we endow it with the induced topology. It is called the {\it maximal spectrum} of $A$. Let
$X$ be a compact Hausdorff space, and let $C(X)$ denote the ring of all real-valued continuous functions on $X$. For each $x \in X$, let $\mathfrak{m}_x$ be the set of all $f \in C(X)$ such that
$f(x) = 0$. The ideal $\mathfrak{m}_x$ is maximal since it is the kernel of the map $C(X) \to \R$ via $f \mapsto f(x)$. Thus, we have a map $\mu : X \to \Max(C(X))$ via $x \mapsto \mathfrak{m}_x$. We aim
to show that $\mu$ is a homeomorphism of $X$ onto $\Max(C(X))$.

\begin{proof}
Let $\mathfrak{m} \in \Max(C(X))$, and define $V = V(\mathfrak{m})$ by
$$V = \{x \in X : f(x) = 0 \text{ for all } f \in \mathfrak{m}\}.$$
Suppose $V$ is empty. Then, for each $x \in X$, there exists $f_x \in \mathfrak{m}$ such that $f_x(x) \neq 0$. Since $f_x$ is continuous, there is an open neighborhood $U_x$ of $x$ in $X$ on which $f_x$
does not vanish. By compactness, a finite number of the neighborhoods, say $U_{x_1},\ldots,U_{x_n}$ cover $X$. Let
$$f = f_{x_1}^2 + \cdots + f_{x_n}^2.$$
Then, $f$ does not vanish anywhere on $X$, hence $1/f$ is a well-defined real-valued function on $X$. Open intervals are a basis of $\R$, so since $0$ is not in the image of $1/f$, it is enough to consider
intervals $(a,b)$, where $0 < a < b$ or $a < b < 0$. Note that $(1/f)^{-1}((a,b)) = f^{-1}((1/b,1/a))$, but this is open by the continuity of $f$, so $1/f$ is also continuous. Hence, $1/f \in C(X)$ and $f$ is
a unit, violating the maximality of $\mathfrak{m}$. Hence, $V$ is non-empty, and there is some $x \in V$. Then, $\mathfrak{m} \subset \mathfrak{m}_x$, so by maximality, $\mathfrak{m} = \mathfrak{m}_x$. 
Hence, $\mu$ is surjective.

Since $X$ is compact and Hausdorff, it is also normal. By the Hausdorff property, with distinct $x,y \in X$, we have disjoint open sets $U_x$ and $U_y$ such that $x \in U_x$ and $y \in U_y$. By normality, there
are disjoint closed sets $V_x$ and $V_y$ containing $U_x$ and $U_y$, respectively. By Urysohn's lemma (since $X$ is normal), there exists a continuous function which is $0$ on $V_x$ and $1$ on $V_y$, and vice-versa.
Thus, $\mathfrak{m}_x$ contains an element that $\mathfrak{m}_y$ does not and vice-versa, so $\mu: x \mapsto \mathfrak{m}_x$ is injective. Compiling results, we've shown that $\mu$ is bijective.

Let $f \in C(X)$. Define
$$U_f = \{x \in X : f(x) \neq 0\} \quad\text{ and }\quad \tilde U_f = \{\mathfrak{m} \in \Max(C(X)) : f \not\in \mathfrak{m}\}.$$
Consider any closed set $C$ of $X$. If $y \not\in C$, then since $\{y\}$ is closed, we have by Urysohn's lemma that there is some $f \in C(X)$ such that $f(y) = 1$ and $f$ is $0$ on $C$. Thus, $C$ is the
intersection of $U_f^c$ for all $f \in C(X)$ such that $f$ is zero on $C$. Thus, any open set is a union of $U_f$, so $\{U_f\}$ forms a basis of $X$.
Also, for any $f \in C(X)$, $X_f$ forms a basis of $\Spec(C(X))$ by Exercise \ref{1.17}. Restricted to $\Max(C(X))$, these open sets have the form
$$\{\mathfrak{m} \in \Max(C(X)) : (f) \subset \mathfrak{m}\}^c = \{\mathfrak{m} \in \Max(C(X)) : (f) \not\subset \mathfrak{m}\} = \{\mathfrak{m} \in \Max(C(X)) : f \not\in \mathfrak{m}\} = \tilde U_f,$$
so $\{\tilde U_f\}$ defines a basis for $\Max(C(X))$ since the topology is inherited from $\Spec(C(X))$.

Thus, if we can show $\mu(U_f) = \tilde U_f$, we will have shown that $\mu$ is an open map and thus a homeomorphism by its already established bijectivity.
If $f(x) \neq 0$, then $f \not\in \mathfrak{m}_x$. Conversely, if $f \not\in \mathfrak{m}_x$, then by definition, $f(x) \neq 0$. Thus, $\mu(U_f) = \tilde U_f$, completing the proof.
\end{proof}

\begin{remark}
This implies the topological space $X$ can be reconstructed from the ring $C(X)$ via $\Max (C(X))$.
\end{remark}

\subsection{Exercise}\label{1.27}

Let $k$ be an algebraically closed field, and let
$$f_a (x_1,\ldots,x_n) = 0$$
be a set of polynomial equations in $n$ variables with coefficients in $k$. The set $X$ of all points $x = (x_1,\ldots,x_n) \in k^n$ which satisfy the equations is an {\it affine algebraic variety}.
Consider the set of all polynomials $g \in k[x_1,\ldots,x_n]$ with the property that $g(x)=0$ for all $x \in X$. This set is an ideal $I(X)$ in the polynomial ring, and is called the {\it ideal of the variety} 
$X$. The quotient ring
$$P(X) = k[x_1,\ldots,x_n] / I(X)$$
is the ring of polynomial functions on $X$ because two polynomials $g,h$ define the same polynomial on $X$ if and only if $g - h$ vanishes at every point of $X$, i.e. if and only if $g-h \in I(X)$.

Let $\xi_i$ be the image of $t_i$ under the projection $k[x_1,\ldots,x_n] \to P(X)$. These are the {\it coordinate functions} on $P(X)$. If $x \in X$, then $\xi_i(x)$ is well-defined and gives the
$i$th coordinate of $x$. Clearly, $P(X)$ is generated as a $k$-algebra by the coordinate functions, and is called the {\it coordinate ring} (or {\it affine algebra}) of $X$.

As in Exercise \ref{1.26}, for each $x \in X$, let $\mathfrak{m}_x$ be the ideal of all $f \in P(X)$ such that $f(x) = 0$; it is a maximal ideal of $P(X)$. Hence, we have a mapping $\mu : X \to \Max(P(X))$
by $x \mapsto \mathfrak{m}_x$. It is easy to show that $\mu$ is injective: if $x \neq y$, we have $\xi_i(x) \neq \xi_i(y)$ for some $i$, so $\xi_i - \xi_i(x) \in \mathfrak{m}_x$, but it is not in $\mathfrak{m}_y$. 
Therefore, $\mathfrak{m}_x \neq \mathfrak{m}_y$. We now show that $\mu$ is surjective; this creates a one-to-one correspondence between $X$ and $\Max(P(X))$, one form of Hilbert's Nullstellensatz.

We'll start by showing that the maximal ideals of $k[x]$ are exactly $(x - a)$ for $a \in k$.
\begin{proof}
Clearly, $(x - a)$ is maximal as it is the kernel of the surjective ring homomorphism $k[x] \to k$ via $f \mapsto f(a)$ since $k$ is algebraically closed. Let $\mathfrak{m}$ be a maximal ideal of $k[x]$. Let
$f \in \mathfrak{m}$. Since the degree of a polynomial is a nonnegative integer, there is a minimal degree $m$ of $\mathfrak{m}$ (it's clear that $\mathfrak{m} \neq 0$). Note that $m > 0$; otherwise $\mathfrak{m}$ contains a unit,
violating maximality. Let $f$ be a polynomial of degree $m$. Then, since $k$ is algebraically closed, it factors as $a_0(x-a_1)\cdots(x-a_m)$ for $a_0 \neq 0$. Redefine $f$ as $a_0^{-1}f$. Since $\mathfrak{m}$ is an ideal 
and $k$ is a subring of $k[x]$, $f = (x-a_1)\cdots(x-a_m)$ is in $\mathfrak{m}$. If $g$ is also a monic degree $m$ polynomial of $\mathfrak{m}$, then if $f-g$ is nonzero, it has degree less than $m$, violating the minimality
of $m$. Hence, $f = g$, so all degree $m$ polynomials in $\mathfrak{m}$ have $a_1,\ldots,a_m$ as roots.

Let $g$ be a polynomial in $\mathfrak{m}$ of minimal degree such that $g(a_1) \neq 0$. Then, it has degree $n > m$ since $f$ is not proportional to $g$ (values at $a_1$). Multiply by some element of $k$ to make it monic; 
it's still in $\mathfrak{m}$. Since $k$ is algebraically closed, $g(x) = (x-b_1)\cdots(x-b_n)$. Now, $f(x)(x-b_{m+1})\cdots(x-b_n)$ is in $\mathfrak{m}$ since $(x-b_{m+1})\cdots(x-b_n) \in k[x]$. Thus,
$$h(x) = f(x)(x-b_{m+1})\cdots(x-b_n) - g(x) = (f(x) - (x-b_1)\cdots(x-b_m))(x-b_{m+1})\cdots(x-b_n).$$
Note that $h(x) \in \mathfrak{m}$. Also, $f(x) - (x-b_1)\cdots(x-b_m) \neq 0$ since $f(a_1) = 0$, but $(x-b_1)\cdots(x-b_m) \neq 0$. However, both are monic of degree $m$, so their difference has
degree less than $m$. Thus, $h(x)$ has degree less than $n$. However, $f(x)(x-b_{m+1})\cdots(x-b_n)$ evaluates to $0$ at $a_1$, but $g(a_1)\neq 0$, so $h(a_1)\neq 0$. This violates the minimality of $n$. Thus,
we have that no such $g$ exists, meaning that $g(a_1) = 0$ for all $g \in \mathfrak{m}$. Hence, $\mathfrak{m} \subset (x - a_1)$, so by maximality, $\mathfrak{m} = (x-a_1)$.
\end{proof}

Now, we show that the maximal ideals of $k[x_1,\ldots,x_n]$ are exactly $(x_1 - a_1,\ldots,x_n - a_n)$ for $a_i \in k$.
\begin{proof}
Consider the ideal $(x_1 - a_1, \ldots, x_n - a_n)$ of $k[x_1,\ldots,x_n]$ for $a_i \in k$. It is the kernel of the surjective ring homomorphism $k[x_1,\ldots,x_n] \to k$ via $f \mapsto f(a_1,\ldots,a_n)$,
so it is a maximal ideal. 

Conversely, suppose $\mathfrak{m}$ is a maximal ideal of $k[x_1,\ldots,x_n]$. For each $i$, we have an injective ring homomorphism $k[x_i] \to k[x_1,\ldots,x_n]$. The contraction of $\mathfrak{m}$ is
a prime ideal of $k[x_i]$. This is contained in some maximal ideal $(x_i - a_i)$ by the previous proof. Note that $\mathfrak{m} \subset (x_i - a_i) + \mathfrak{m}$; since the contraction of $\mathfrak{m}$ is 
contained in $(x_i - a_i)$, the contraction of the sum is also contained in $(x_i - a_i)$. Thus, the sum cannot be the whole ring, so by the maximality of $\mathfrak{m}$, we have $(x_i - a_i) \subset \mathfrak{m}$.
Doing this for each $i$, we have that $(x_i - a_i)$ is a subset of $\mathfrak{m}$ for each $i$. Thus, $(x_1 - a_1,\ldots,x_n - a_n) \subset \mathfrak{m}$, so by the maximality of both sets,
$\mathfrak{m} = (x_1 - a_1,\ldots,x_n - a_n)$.
\end{proof}

Now, we can prove the surjectivity of $\mu$.

\begin{proof}
Suppose $\mathfrak{m} \in \Max(P(X))$. Using the ring homomorphism $\phi : k[x_1,\ldots,x_n] \to P(X)$, we have that $\phi^{-1}(\mathfrak{m})$ is a prime ideal of $k[x_1,\ldots,x_n]$. Let $\mathfrak{M}$
be a maximal ideal of $k[x_1,\ldots,x_n]$ containing $\phi^{-1}(\mathfrak{m})$. Since $\phi$ is surjective, $\phi(\mathfrak{M})$ is an ideal of $P(X)$ containing $\mathfrak{m}$. By maximality, $\phi(\mathfrak{M}) = \mathfrak{m}$
or $\phi(\mathfrak{M}) = P(X)$. 

In the first case, $\mathfrak{M} \subset \phi^{-1}(\mathfrak{m})$, so $\mathfrak{M} = \phi^{-1}(\mathfrak{m})$, and $\phi^{-1}(\mathfrak{m})$ is a maximal ideal of $k[x_1,\ldots,x_n]$. In the second case,
since $P(X) = k[x_1,\ldots,x_n] / I(X)$, we have that $I(X) + \mathfrak{M} = k[x_1,\ldots,x_2]$. In particular, since $\mathfrak{M}$ is maximal, $\mathfrak{M}$ does not contain $I(X)$; else, 
$k[x_1,\ldots,x_2] = I(X) + \mathfrak{M} = \mathfrak{M}$. But, $I(X) \subset \phi^{-1}(\mathfrak{m})$, violating the fact that $\mathfrak{M}$ contains $\phi^{-1}(\mathfrak{m})$. Hence, we have
that $\phi^{-1}(\mathfrak{m})$ is a maximal ideal of $k[x_1,\ldots,x_n]$.

By our previous proof, we then have that $\phi^{-1}(\mathfrak{m}) = (x_1 - a_1, \ldots, x_n - a_n)$ for $a_i \in k$. Since $I(X) \subset \phi^{-1}(\mathfrak{m})$, we must have $a = (a_1,\ldots,a_n) \in X$. By
the surjectivity of the projection map $\phi$, $\phi((x_1 - a_1, \ldots, x_n - a_n)) = \mathfrak{m}$. Since for any $f \in \phi((x_1 - a_1, \ldots, x_n - a_n))$, $f(a) = 0$, we have that $\mathfrak{m}_a \subset \mathfrak{m}$.
By maximality, $\mathfrak{m} = \mathfrak{m}_a$, proving the surjectivity of $\mu$.
\end{proof}

\subsection{Exercise}\label{1.28}

Let $f_1, \ldots, f_m$ be elements of $k[x_1,\ldots,x_n]$. They determine a {\it polynomial mapping} $\phi : k^n \to k^m$; if $x \in k^n$, the cooordinates of $\phi(x)$ are $f_1(x),\ldots,f_m(x)$.
Let $X$, $Y$ be affine algebraic varieties in $k^n$, $k^m$, respectively. A mapping $\phi : X \to Y$ is said to be {\it regular} if $\phi$ is the restriction to $X$ of a polynomial mapping from
$k^n$ to $k^m$. If $\eta$ is a polynomial function on $Y$, then $\eta \circ \phi$ is a polynomial function on $X$. Hence $\phi$ induces a $k$-algebra homomorphism $P(Y) \to P(X)$, namely $\eta \mapsto \eta \circ \phi$.
We will show that this yields a one-to-one correspondence between the regular mappings $X \to Y$ and the $k$-algebra homomorphisms $P(Y) \to P(X)$.

\begin{proof}
Let $\phi : X \to Y$ be a regular mapping. We'll start by showing that the $k$-homomorphism $\eta \mapsto \eta \circ \phi$ is unique to $\phi$. Suppose $\psi : X \to Y$ is a regular function and
$\eta \circ \psi = \eta \circ \phi$ for all $\eta \in P(Y)$. Consider the coordinate functions $\xi_i \in P(Y)$. Then, $\xi_i \circ \psi = \xi_i \circ \phi$ for all $i$, so $\psi = \phi$.

Now, consider an arbitrary $k$-algebra homomorphism $\theta: P(Y) \to P(X)$, or
$$\theta: k[y_1,\ldots,y_m]/I(Y) \to k[x_1,\ldots,x_n]/I(X).$$
Since $\xi_i$ generate $P(Y)$ as a $k$-algebra, $\theta$ is defined entirely by $\theta(\xi_i)$. Define $\theta_1$ to be the map
$$\theta_1: k[y_1,\ldots,y_m] \to k[x_1,\ldots,x_n]/I(X)$$
that sends $y_i \mapsto \theta(\xi_i)$. Then, $\theta_1$ is a $k$-algebra homomorphism. Likewise, there are $\Xi_i$ that generate $P(X)$. Hence, 
$\theta_1(y_i) = \sum_{\alpha \in I_i}c_{i,\alpha}\Xi^{\alpha}$ for $c_{i,\alpha} \in k$. Then, we can define $\psi$ to be the unique $k$-algebra homomorphism
$$\psi : k[y_1,\ldots,y_m] \to k[x_1,\ldots,x_n]$$
defined by $\psi(y_i) = \sum_{\alpha \in I_i}c_{i,\alpha}x^{\alpha}$. Then, if the canonical projection maps defining the varieties are $\pi_Y:k[y_1,\ldots,y_m] \to k[y_1,\ldots,y_m] /I(Y)$ and 
$\pi_X : k[x_1,\ldots,x_n] \to k[x_1,\ldots,x_n]/I(X)$, we have $\pi_X \circ \psi = \theta \circ \pi_Y$ by our definition of $\psi$. In particular, $\pi_X(\psi(I(Y))) = \theta(0) = 0$, so $\psi : I(Y) \to I(X)$. 

For each $y_i$, $\psi(y_i)$ is a polynomial of $k[x_1,\ldots,x_n]$. Hence, $(\psi(y_1)(x),\ldots,\psi(y_m)(x))$ defines a map $k[x_1,\ldots,x_n] \to k[y_1,\ldots,y_m]$. Since
$\psi : I(Y) \to I(X)$, this defines a map $X \to Y$. This is regular as it is the restriction of a $k$-homomorphism $k[x_1,\ldots,x_n] \to k[y_1,\ldots,y_m]$.

The regular functions $X \to Y$ are one-to-one with the $k$-homomorphisms $P(Y) \to P(X)$.
\end{proof}
